\documentclass[11pt]{article}
\usepackage[margin=1in]{geometry}
\usepackage{amsmath,amssymb,amsthm}
\usepackage[hidelinks]{hyperref}

\newcommand{\E}{\mathbb E}
\newcommand{\Var}{\operatorname{Var}}
\newcommand{\CV}{\operatorname{CV}}
\newtheorem{theorem}{Theorem}
\newtheorem{lemma}[theorem]{Lemma}
\newtheorem{proposition}[theorem]{Proposition}
\newtheorem{corollary}[theorem]{Corollary}
\newtheorem{conjecture}[theorem]{Conjecture}
\theoremstyle{remark}
\newtheorem{remark}[theorem]{Remark}

\title{A Fisher-Information Reduction of the Cyclic Phase-Type Moment Hierarchy}
\author{}
\date{}

\begin{document}
\maketitle

\section{The reduction}

Let \(T\) be the absorption time of a phase-type chain with \(s\) transient
states, and write
\[
  m_j=\E T^j.
\]
For \(r\geq0\), let \(T^{[r]}\) have the \(r\)-fold size-biased density
\[
  g_r(t)=\frac{t^r f(t)}{m_r},\qquad t>0.
\]
Then
\[
  \E T^{[r]}=\frac{m_{r+1}}{m_r},\qquad
  \CV^2(T^{[r]})
  =\frac{m_rm_{r+2}}{m_{r+1}^2}-1.
\]
Consequently the desired hierarchy at \(k=r+1\),
\[
  m_{k-1}m_{k+1}\geq
  \frac{s+k}{s+k-1}m_k^2,
\]
is equivalent to
\begin{equation}\label{eq:cv-target}
  \CV^2(T^{[r]})\geq\frac1{s+r}.
\end{equation}

The following analytic condition is sufficient for all the inequalities
\eqref{eq:cv-target} simultaneously.

\begin{proposition}[Log-convex density quotient implies the hierarchy]
\label{prop:fisher}
Suppose
\begin{equation}\label{eq:qdef}
  q(t):=\frac{f(t)}{t^{s-1}}
\end{equation}
is log-convex on \((0,\infty)\).  Then \eqref{eq:cv-target} holds for every
\(r\geq0\).
\end{proposition}

\begin{proof}
Put \(d=s+r\).  Apart from its normalizing constant, the density of
\(T^{[r]}\) is
\[
  g(t)=t^{d-1}q(t).
\]
Let
\[
  \ell(t)=\frac{q'(t)}{q(t)}.
\]
The score for the log-scale family \(g_a(t)=e^a g(e^a t)\) at \(a=0\) is
\[
  S(t)=1+t\frac{g'(t)}{g(t)}=d+t\ell(t).
\]
Integration by parts gives
\begin{equation}\label{eq:scoremeans}
  \E S=0,\qquad \E[t\ell(t)]=-d.
\end{equation}
Log-convexity of \(q\) says
\[
  \ell'(t)\geq0,
  \qquad\text{equivalently}\qquad
  \ell(t)^2\leq\frac{q''(t)}{q(t)}.
\]
Two integrations by parts therefore give
\[
  \E[(t\ell(t))^2]
  \leq
  \frac{\int_0^\infty t^{d+1}q''(t)\,dt}
       {\int_0^\infty t^{d-1}q(t)\,dt}
  =d(d+1).
\]
It follows from \eqref{eq:scoremeans} that the scale Fisher information is
\[
  I=\E[S^2]
   =\E[(t\ell(t))^2]-d^2
   \leq d.
\]
Finally, another integration by parts yields
\[
  \E[(T^{[r]}-\E T^{[r]})S(T^{[r]})]
  =-\E T^{[r]}.
\]
Cauchy--Schwarz now gives
\[
  \Var(T^{[r]})I\geq(\E T^{[r]})^2.
\]
Since \(I\leq d=s+r\), this is \eqref{eq:cv-target}.
\end{proof}

All boundary terms in this calculation vanish for a phase-type density,
including when the representation permits immediate absorption or the density
has a lower-than-\(s-1\) order zero at the origin.

\section{Uniformization and the one missing matrix lemma}

Write a phase-type density as
\[
  f(t)=\alpha e^{Qt}b,
  \qquad b=-Q\mathbf1.
\]
Choose a uniformization rate
\(\lambda\geq\max_i(-Q_{ii})\), and put
\[
  P=I+\lambda^{-1}Q\geq0,
  \qquad A=\lambda P\geq0.
\]
Then
\[
  f(t)=e^{-\lambda t}\alpha e^{tA}b.
\]
Multiplication by \(e^{-\lambda t}\) does not change the second derivative of
the logarithm.  Thus Proposition~\ref{prop:fisher} would follow from the next
standalone statement.

\begin{conjecture}[Dimension-corrected matrix-exponential log-convexity]
\label{conj:matrix}
Let \(A\) be an entrywise nonnegative \(s\times s\) matrix and let
\(u,v\geq0\).  If
\[
  F(t)=u^T e^{tA}v
\]
is not identically zero, then
\begin{equation}\label{eq:matrix-logconvex}
  t^{1-s}F(t)\quad\text{is log-convex on }(0,\infty).
\end{equation}
Equivalently,
\begin{equation}\label{eq:matrix-curvature}
  t^2\bigl(F(t)F''(t)-F'(t)^2\bigr)+(s-1)F(t)^2\geq0.
\end{equation}
\end{conjecture}

There is an equivalent local differential formulation which removes the two
endpoint vectors.  It appears well suited to an induction on the directed
graph.

\begin{proposition}[Growth-rate formulation]\label{prop:growth-rates}
Let $h(t)=e^{tA}v$, initially with $v>0$, and put
\[
 r_i(t)=\frac{h_i'(t)}{h_i(t)},\qquad
 c_{ij}(t)=A_{ij}\frac{h_j(t)}{h_i(t)}.
\]
Then
\begin{equation}\label{eq:rate-system}
 r_i'(t)=\sum_j c_{ij}(t)\bigl(r_j(t)-r_i(t)\bigr).
\end{equation}
Conjecture~\ref{conj:matrix} is equivalent to the componentwise estimates
\begin{equation}\label{eq:rate-target}
 r_i'(t)\geq-\frac{s-1}{t^2}
 \qquad(1\leq i\leq s).
\end{equation}
\end{proposition}

\begin{proof}
Differentiating $r_i=h_i'/h_i$, using $h'=Ah$, gives
\[
 r_i'
 =\frac{\sum_jA_{ij}h_jr_j}{h_i}-r_i^2
 =\sum_jc_{ij}(r_j-r_i),
\]
because $\sum_jc_{ij}=r_i$.  Taking $u=e_i$ in the matrix conjecture
shows that \eqref{eq:rate-target} is necessary.  Conversely, for arbitrary
$u\geq0$, set
\[
 \pi_i(t)=\frac{u_i h_i(t)}{u^Th(t)}.
\]
A direct differentiation, or the usual log-sum identity, gives
\[
 (\log(u^Th))''
 =\sum_i\pi_i r_i'
   +\sum_i\pi_i\left(r_i-\sum_j\pi_jr_j\right)^2.
\]
Thus \eqref{eq:rate-target} implies
$(\log(u^Th))''\geq-(s-1)/t^2$.  Vectors with zero entries follow by
approximation.
\end{proof}

The coefficients $c_{ij}$ in \eqref{eq:rate-system} are nonnegative and the
matrix $(c_{ij})$ is diagonally similar to $A$.  Moreover, along every
directed cycle their product is independent of $t$, since the factors
$h_j/h_i$ telescope.  Hence a negative $r_i'$ can only be produced by flow
from $i$ toward vertices with smaller growth rate, while every closed cycle
has zero total logarithmic change.  The remaining problem is to turn this
finite descending-chain structure into the sharp loss $(s-1)/t^2$.

\begin{theorem}[Conditional all-\(k\) cyclic theorem]
Conjecture~\ref{conj:matrix} implies the phase-type moment hierarchy for every
number of states \(s\) and every \(k\geq1\), including cyclic chains.
\end{theorem}

\begin{proof}
Apply Conjecture~\ref{conj:matrix} with
\(u^T=\alpha\) and \(v=b\).  Uniformization shows that the quotient
\(f(t)/t^{s-1}\) is log-convex.  Proposition~\ref{prop:fisher} then proves all
the adjacent-moment inequalities.
\end{proof}

This isolates the cyclic problem in a form that no longer mentions absorption
times, size biasing, or moments.  In fact, if the moment hierarchy is required
for every exponential tilt of every phase-type law, the matrix statement is
also necessary.  We record the precise converse next.

\begin{lemma}[Laplace concentration]\label{lem:laplace-concentration}
Let $H$ be positive and $C^4$ near $t_0>0$, with enough exponential
decay that the following integrals exist.  Put
\[
 J_{r,j}=\int_0^\infty
          t^{r+j}e^{-rt/t_0}H(t)\,dt
 \qquad(j=0,1,2).
\]
Then, as $r\to\infty$,
\begin{equation}\label{eq:laplace-expansion}
 \frac{J_{r,0}J_{r,2}}{J_{r,1}^2}-1
 =\frac1r+
   \frac{t_0^2(\log H)''(t_0)-1}{r^2}
   +o(r^{-2}).
\end{equation}
\end{lemma}

\begin{proof}
The probability measure proportional to
$t^re^{-rt/t_0}H(t)\,dt$ is a smooth tilt of a
$\operatorname{Gamma}(r+1,r/t_0)$ law.  Taylor-expand
$\log H(t)$ at $t_0$, use the centered Gamma moments through order four,
and expand its squared coefficient of variation.  The constant and linear
Taylor terms only change normalization and scale; the quadratic term gives
$t_0^2(\log H)''(t_0)/r^2$.  The unperturbed Gamma law contributes
$1/(r+1)=r^{-1}-r^{-2}+O(r^{-3})$, which proves
\eqref{eq:laplace-expansion}.  Standard Laplace remainder estimates make the
expansion uniform on a fixed neighborhood of $t_0$.
\end{proof}

\begin{proposition}[The matrix lemma is forced by the tilted hierarchy]
\label{prop:converse}
Suppose the phase-type hierarchy holds, for a fixed state bound $s$, for
every $k$ and every $s$-state phase-type law.  Then
Conjecture~\ref{conj:matrix} holds for every nonnegative $s\times s$ matrix
$A$ and $u,v\geq0$.
\end{proposition}

\begin{proof}
It is enough first to take $u,v>0$; the general case follows by
approximation.  Choose $\lambda>\rho(A)$, put
\[
 d=(\lambda I-A)^{-1}v>0,\qquad D=\operatorname{diag}(d),
\]
and define
\[
 Q=D^{-1}AD-\lambda I,\qquad b=D^{-1}v.
\]
Then $Q$ is a phase-type subgenerator and $-Q\mathbf1=b$.  After multiplying
$u$ by an irrelevant positive constant, take
$\alpha=u^TD$ to be a probability vector.  Its absorption density
is
\[
 \alpha e^{tQ}b=e^{-\lambda t}u^Te^{tA}v.
\]
For every $\theta>0$, exponentially killing this chain at rate $\theta$
and conditioning on absorption through the original exit $b$ gives another
$s$-state phase-type law.  Its density is proportional to
\[
 e^{-\theta t}\alpha e^{tQ}b
 =e^{-(\lambda+\theta)t}F(t),
 \qquad F(t)=u^Te^{tA}v.
\]
Indeed this is the usual Doob transform with
$h=(\theta I-Q)^{-1}b$.

Fix $t_0>0$.  For large $r$, choose
$\theta=r/t_0-\lambda>0$, and apply the hierarchy at $k=r+1$.
The squared coefficient of variation of the $r$-fold size-biased law is
the left side of \eqref{eq:laplace-expansion} with $H=F$, whereas the
hierarchy says that it is at least $1/(s+r)$.  Hence
\[
 \frac1r+\frac{t_0^2(\log F)''(t_0)-1}{r^2}+o(r^{-2})
 \geq
 \frac1r-\frac{s}{r^2}+O(r^{-3}).
\]
Multiplying by $r^2$ and letting $r\to\infty$ gives
\[
 t_0^2(\log F)''(t_0)+(s-1)\geq0.
\]
Since $t_0$ was arbitrary, this is
\eqref{eq:matrix-curvature}.
\end{proof}

\section{The complete two-state cyclic case}

The matrix conjecture, and hence the full hierarchy, can be proved
unconditionally when \(s=2\).

\begin{proposition}\label{prop:two-by-two}
If \(A\) is a nonnegative \(2\times2\) matrix and \(u,v\geq0\), then
\[
  \frac{u^Te^{tA}v}{t}
\]
is log-convex on \((0,\infty)\).
\end{proposition}

\begin{proof}
Write
\[
 A=\begin{pmatrix}a&b\\c&d\end{pmatrix},\qquad
 \mu=\frac{a+d}{2},\qquad
 \delta=\sqrt{\left(\frac{a-d}{2}\right)^2+bc}.
\]
When \(\delta>0\),
\[
 e^{tA}=e^{\mu t}\left(
   \cosh(\delta t)I+
   \frac{\sinh(\delta t)}{\delta}(A-\mu I)\right).
\]
The off-diagonal entries divided by \(t\) are nonnegative constant
multiples of
\[
 e^{\mu t}\frac{\sinh(\delta t)}{\delta t}.
\]
This is log-convex because
\[
 \frac{d^2}{dt^2}\log\frac{\sinh(\delta t)}{\delta t}
 =\frac1{t^2}-\delta^2\operatorname{csch}^2(\delta t)\geq0.
\]
For a diagonal entry, put
\[
 h=\frac{a-d}{2\delta},\qquad |h|\leq1.
\]
Apart from the common factor \(e^{\mu t}\), that entry is either
\[
 \cosh(\delta t)+h\sinh(\delta t)
 =\frac{1+h}{2}e^{\delta t}+\frac{1-h}{2}e^{-\delta t}
\]
or the same expression with \(h\) replaced by \(-h\).  It is therefore a
positive sum of exponentials and is log-convex.  Division by \(t\) preserves
log-convexity because \(-\log t\) is convex.  The case \(\delta=0\) follows by
continuity.  Finally, \(u^Te^{tA}v/t\) is a positive sum of these four
log-convex entries, and positive sums of log-convex functions are
log-convex.
\end{proof}

\begin{corollary}[All moments for two cyclic states]
Every phase-type law with a representation of size at most two satisfies
\[
 m_{k-1}m_{k+1}\geq\frac{k+2}{k+1}m_k^2
 \qquad(k\geq1).
\]
\end{corollary}

\begin{proof}
Apply Proposition~\ref{prop:two-by-two} after uniformization and then use
Proposition~\ref{prop:fisher}.
\end{proof}

\section{Closure under block-triangular composition}

The two-state result extends immediately to chains made from one- and
two-state strongly connected components.  The key closure principle is useful
more generally.

\begin{lemma}[Dimension-additive convolution]\label{lem:convolution}
Let $p,q\geq1$.  Suppose $f,g\geq0$, and suppose
\[
 t^{1-p}f(t)\quad\hbox{and}\quad t^{1-q}g(t)
\]
are log-convex.  Then, for
\[
 (f*g)(t)=\int_0^t f(t-x)g(x)\,dx,
\]
the function $t^{1-p-q}(f*g)(t)$ is log-convex.
\end{lemma}

\begin{proof}
Write $f(t)=t^{p-1}\phi(t)$ and
$g(t)=t^{q-1}\psi(t)$.  The substitution $x=ty$ gives
\[
 \frac{(f*g)(t)}{t^{p+q-1}}
 =\int_0^1(1-y)^{p-1}y^{q-1}
   \phi(t(1-y))\psi(ty)\,dy.
\]
For fixed $y$, the integrand is log-convex as a function of $t$.
Positive integrals of log-convex functions are log-convex by H\"older's
inequality.
\end{proof}

\begin{proposition}[Block-triangular closure]\label{prop:block-closure}
Suppose Conjecture~\ref{conj:matrix} holds in dimensions $p$ and $q$.
Then it holds for every nonnegative block upper-triangular matrix
\[
 A=\begin{pmatrix}B&C\\0&E\end{pmatrix},
 \qquad B\in\mathbb R^{p\times p},\quad
 E\in\mathbb R^{q\times q}.
\]
\end{proposition}

\begin{proof}
It is enough to prove the assertion entrywise, since a positive sum of
log-convex functions is log-convex.  Entries in a diagonal block satisfy the
claim by the corresponding lower-dimensional conjecture, followed by
multiplication by $t^{-q}$ or $t^{-p}$.  For the upper-right block, Duhamel's
formula gives
\[
 (e^{tA})_{ij}
 =\sum_{a,b}C_{ab}\int_0^t
   (e^{(t-x)B})_{ia}(e^{xE})_{bj}\,dx.
\]
Every summand has the required dimension-$p+q$ correction by
Lemma~\ref{lem:convolution}; summing preserves log-convexity.
\end{proof}

\begin{corollary}[Components of size at most two]
If every strongly connected component of the directed graph of a nonnegative
$s\times s$ matrix $A$ has at most two vertices, then
$t^{1-s}u^Te^{tA}v$ is log-convex for all $u,v\geq0$.  Consequently the
full moment hierarchy holds for every phase-type chain whose transient
strongly connected components have size at most two.
\end{corollary}

\begin{proof}
Permute the strongly connected components into block upper-triangular order
and repeatedly apply Proposition~\ref{prop:block-closure}.  Components of size
one are immediate, and components of size two are covered by
Proposition~\ref{prop:two-by-two}.  The phase-type conclusion follows from
Proposition~\ref{prop:fisher}.
\end{proof}

\section{An exact reduction of the three-state problem}

There is a useful way to pass from an arbitrary entry of a positive
three-dimensional semigroup to an absorption density.  This removes the two
endpoint vectors and makes the order-three canonical representation directly
applicable.

\begin{lemma}[A matrix entry is a phase-type density]\label{lem:entry-to-ph}
Let $A$ be an irreducible nonnegative $s\times s$ matrix and fix $i,j$.
There is a transient Markov generator $T$ with $s$ transient states, an exit
vector $b=-T\mathbf1$ supported only at $j$, and constants $c,\lambda>0$ such
that
\[
 (e^{tA})_{ij}=c e^{\lambda t}e_i^Te^{tT}b.
\]
Thus the logarithmic curvatures of the two sides are equal.
\end{lemma}

\begin{proof}
Choose $\lambda>\rho(A)$ and put
\[
 h=(\lambda I-A)^{-1}e_j
   =\sum_{n\geq0}\lambda^{-n-1}A^ne_j.
\]
Irreducibility gives $h>0$.  With $H=\operatorname{diag}(h)$, set
\[
 T=H^{-1}(A-\lambda I)H.
\]
Its off-diagonal entries are nonnegative, and
\[
 T\mathbf1=-H^{-1}e_j,
 \qquad b=H^{-1}e_j.
\]
The killing state is reachable from every transient state, so $T$ is
transient.  Finally,
\[
 e_i^Te^{tT}b
 =h_i^{-1}e^{-\lambda t}(e^{tA})_{ij}.
\]
Take $c=h_i$.  Multiplication by a positive constant and by an exponential
does not change logarithmic curvature.
\end{proof}

We now use the He--Zhang unicyclic representation theorem: every phase-type
law of order three has a distribution-equivalent Markovian representation
with subgenerator
\begin{equation}\label{eq:canonical-three}
 T_c=\begin{pmatrix}
 -x_1&0&\eta x_1\\
 x_2&-x_2&0\\
 0&x_3&-x_3
 \end{pmatrix},
 \qquad
 x_1\geq x_2\geq x_3>0,\quad 0\leq\eta<1.
\end{equation}
Absorption occurs only from state $1$, at rate $(1-\eta)x_1$.  The canonical
initial distribution need not be deterministic, but this causes no loss:
after division by $t^2$ its density is a positive sum of the three
deterministic-start densities.

Put
\[
 D(z)=(z+x_1)(z+x_2)(z+x_3)-\eta x_1x_2x_3.
\]
A direct resolvent calculation gives their Laplace transforms:
\begin{equation}\label{eq:canonical-laplace}
 \begin{aligned}
 L_1(z)&=\frac{(1-\eta)x_1(z+x_2)(z+x_3)}{D(z)},\\
 L_2(z)&=\frac{(1-\eta)x_1x_2(z+x_3)}{D(z)},\\
 L_3(z)&=\frac{(1-\eta)x_1x_2x_3}{D(z)}.
 \end{aligned}
\end{equation}
Consequently the whole $3\times3$ matrix problem has been reduced to the
following scalar statement.

\begin{conjecture}[Canonical cubic lemma]\label{conj:canonical-cubic}
For the parameters in \eqref{eq:canonical-three}, let $f_r$ be the inverse
Laplace transform of $L_r$ in \eqref{eq:canonical-laplace}.  Then
\[
 t^{-2}f_r(t)\quad\hbox{is log-convex on $(0,\infty)$}
 \qquad(r=1,2,3).
\]
\end{conjecture}

\begin{proposition}[The cubic lemma is exactly sufficient in dimension three]
\label{prop:three-reduction}
Conjecture~\ref{conj:canonical-cubic} implies
Conjecture~\ref{conj:matrix} for every nonnegative $3\times3$ matrix.
\end{proposition}

\begin{proof}
It is enough to treat matrix entries.  Reducible matrices are already covered
by Proposition~\ref{prop:block-closure}.  For an irreducible matrix apply
Lemma~\ref{lem:entry-to-ph}, then replace the resulting order-three
phase-type law by its He--Zhang representation.  Its density is a positive
sum of $f_1,f_2,f_3$.  Positive sums preserve log-convexity, and the constant
and exponential factors in Lemma~\ref{lem:entry-to-ph} are harmless.
\end{proof}

The remaining lemma can also be written as a geometric-renewal problem.  Let
\[
 g_r(t)=x_re^{-x_rt},\qquad C=g_1*g_2*g_3,
 \qquad R=\sum_{n\geq0}\eta^n C^{*n},
\]
where $C^{*0}$ is a unit mass at zero.  Then
\begin{equation}\label{eq:canonical-renewal}
 f_1=(1-\eta)g_1*R,\qquad
 f_2=(1-\eta)g_2*g_1*R,\qquad
 f_3=(1-\eta)C*R.
\end{equation}
For $\eta=0$ the assertion follows immediately from
Lemma~\ref{lem:convolution}: these are convolutions of respectively one, two,
and three exponential densities, with extra factors $t^{-2}$, $t^{-1}$, and
$1$.  What remains is therefore precisely the effect of geometric feedback.

There is also a complete coefficientwise proof when the three rates are
equal.  This supplies a nontrivial cyclic subcase and identifies the sharp
root-of-unity calculation behind the directed three-cycle.

\begin{proposition}[Equal-rate canonical cycle]\label{prop:equal-rate-cycle}
Conjecture~\ref{conj:canonical-cubic} holds when
$x_1=x_2=x_3$.
\end{proposition}

\begin{proof}
Scale the common rate to one.  Apart from a positive constant and the common
factor $e^{-t}$, the three densities in \eqref{eq:canonical-renewal} are
\[
 H_r(t)=\sum_{n\geq0}\eta^n\frac{t^{3n+r}}{(3n+r)!},
 \qquad r=0,1,2.
\]
The exponential factor has no effect on logarithmic curvature.  In the
coefficient of $\eta^N$ in
\[
 t^2(H_rH_r''-(H_r')^2)+2H_r^2,
\]
all terms have degree $L=3N+2r$.  After multiplication by the positive factor
$L!$, that coefficient is
\[
 \sum_{\substack{0\leq p\leq L\\p\equiv r\pmod3}}
 \binom Lp\left(2-\frac L2+\frac{(L-2p)^2}{2}\right).
\]
Let $\varepsilon=(-1)^N$.  A roots-of-unity filter, applied first to
$(1+x)^L$ and then to
\[
 (L-2x\partial_x)^2(1+x)^L
 =L(1+x)^{L-2}\bigl(L(1-x)^2+4x\bigr),
\]
evaluates the last display as
\begin{equation}\label{eq:equal-rate-coefficient}
 \frac{2^{L+1}+\varepsilon(-3L^2+3L+4)}{3}.
\end{equation}
This is nonnegative.  For odd $N$ it is immediate.  For even $N$, the only
borderline small case is $L=4$, where equality holds; for $L\geq5$ the
elementary inequality
$2^{L+1}\geq3L^2-3L-4$ follows by induction.  The remaining cases
$L=0,2$ are positive by inspection.  Hence every coefficient of the
curvature expression is nonnegative for $0\leq\eta<1$.

The same calculation gives a useful strengthening for the middle start.
For $r=1$, replace the correction $2H_1^2$ by $H_1^2$, as required for
log-convexity of $t^{-1}H_1$.  The filtered coefficient becomes
\begin{equation}\label{eq:equal-rate-middle-strong}
 \frac{2^L+(-1)^N(-3L^2+3L+2)}{3},
 \qquad L=3N+2.
\end{equation}
For odd $N$ it is positive.  For even $N$, equality holds at $L=2$, and
$2^L\geq3L^2-3L-2$ for $L\geq8$.  Hence $t^{-1}H_1(t)$ is log-convex in
the equal-rate case.
\end{proof}

At this stage of the direct canonical route, the unequal-rate version of
Conjecture~\ref{conj:canonical-cubic} is the only unproved step.  Formula
\eqref{eq:canonical-renewal} suggests proving that unequal exponential rates
can only increase the curvature in the equal-rate certificate, but a valid
polarization or sum-of-squares argument for that assertion is not yet known.

For further work it is useful to remove the negative diagonal without losing
the special graph.  Multiplication by $e^{x_1t}$ does not change logarithmic
curvature.  Put
\[
 B=T_c+x_1I=
 \begin{pmatrix}
 0&0&c\\ x_2&a&0\\0&x_3&b
 \end{pmatrix},\qquad
 a=x_1-x_2,\qquad b=x_1-x_3,\qquad c=\eta x_1,
\]
so that $0\leq a\leq b$.  If $h=e^{tB}e_1$, define
\[
 z=c\frac{h_3}{h_1},\qquad
 p=x_2\frac{h_1}{h_2},\qquad
 q=x_3\frac{h_2}{h_3}.
\]
The three component growth rates are
\[
 r_1=z,\qquad r_2=a+p,\qquad r_3=b+q,
\]
and satisfy the closed cyclic system
\begin{equation}\label{eq:canonical-rate-system}
 z'=z(b+q-z),\qquad
 p'=p(z-a-p),\qquad
 q'=q(a+p-b-q),
 \qquad zpq=cx_2x_3.
\end{equation}
Thus the remaining lemma is equivalent to the three elementary-looking
inequalities
\[
 t^2z(z-b-q)\leq2,\qquad
 t^2p(a+p-z)\leq2,\qquad
 t^2q(b+q-a-p)\leq2.
\]
The initial limits are respectively $0,1,2$.  Extensive numerical tests find
no violation; in fact the middle expression appears bounded by $1$.  This
Riccati formulation uses only five nonnegative scalar parameters and the
ordering $a\leq b$, and is currently the most compact form of the missing
unequal-rate argument.

The stronger middle bound also removes the third inequality altogether.
More precisely, it is enough to prove
\begin{equation}\label{eq:two-riccati-targets}
 t^2z(z-b-q)\leq2,
 \qquad
 t^2p(a+p-z)\leq1.
\end{equation}
Indeed, the second inequality says exactly that $t^{-1}h_2(t)$ is
log-convex.  Since
\[
 \bigl(e^{-bt}h_3(t)\bigr)'=x_3e^{-bt}h_2(t),
 \qquad h_3(0)=0,
\]
we have
\[
 e^{-bt}h_3(t)=x_3\int_0^t e^{-bs}h_2(s)\,ds.
\]
Apply Lemma~\ref{lem:convolution} with $p=2$, $q=1$, and with the
second factor equal to the constant function $1$.  It follows that
$t^{-2}e^{-bt}h_3(t)$, and hence $t^{-2}h_3(t)$, is log-convex.  Thus only
the first inequality and the sharper middle inequality in
\eqref{eq:two-riccati-targets} are independent targets.

Two of the starts can in fact be completed whenever the canonical cubic has
only real roots.  Let
\[
 P(s)=s(s-a)(s-b)-Q,
 \qquad Q=cx_2x_3,
\]
and suppose its roots are
$\lambda_1\leq\lambda_2\leq\lambda_3$.  The Laplace transforms of the
shifted second and third components are
\[
 \mathcal L h_2(s)=\frac{x_2(s-b)}{P(s)},
 \qquad
 \mathcal L h_3(s)=\frac{x_2x_3}{P(s)}.
\]
Since $P(b)=-Q<0$ and $P(s)\to\infty$ as $s\to\infty$, one has
$\lambda_1\leq b<\lambda_3$.  With
\[
 \theta=\frac{b-\lambda_1}{\lambda_3-\lambda_1}\in[0,1],
\]
partial fractions give
\begin{equation}\label{eq:middle-divided-difference}
 \frac{s-b}{P(s)}
 =\frac{\theta}{(s-\lambda_1)(s-\lambda_2)}
  +\frac{1-\theta}{(s-\lambda_2)(s-\lambda_3)}.
\end{equation}
The inverse transform of each summand, divided by $t$, is a positive
mixture of exponentials, since
\[
 \frac{e^{ut}-e^{vt}}{u-v}
 =t\int_0^1e^{t((1-y)v+yu)}\,dy.
\]
Consequently $t^{-1}h_2(t)$ is log-convex, proving the sharper middle
inequality in \eqref{eq:two-riccati-targets}.  Likewise, the inverse transform
of $1/P(s)$ is the second divided difference of the exponential:
\[
 h_3(t)=x_2x_3t^2
 \int_{\substack{y_i\geq0\\y_1+y_2+y_3=1}}
 e^{t(y_1\lambda_1+y_2\lambda_2+y_3\lambda_3)}\,d\sigma(y).
\]
Thus $t^{-2}h_3(t)$ is log-convex directly.  Repeated roots follow by
continuity.  In the real-spectrum region, therefore, only the first canonical
start remains open; the genuinely cyclic complex-pair region still requires
the first and middle bounds in \eqref{eq:two-riccati-targets}.

The remaining real-spectrum first start also has a useful positive-mixture
form.  Since
\[
 \frac{(s-a)(s-b)}{P(s)}=\frac1s+\frac{Q}{sP(s)},
\]
the Hermite--Genocchi formula gives
\begin{equation}\label{eq:real-first-simplex}
 h_1(t)=1+w(t),\qquad
 w(t)=Qt^3M(t),\qquad
 M(t)=\int_{\Delta_3}e^{tX(y)}\,dy,
\end{equation}
where
\[
 X(y)=y_1\lambda_1+y_2\lambda_2+y_3\lambda_3,
 \qquad
 \Delta_3=\{y_i\geq0:y_1+y_2+y_3\leq1\}.
\]
Let $\nu_t$ be the exponential tilt of normalized Lebesgue measure on
$\Delta_3$ by $e^{tX}$, and put
\[
 \theta=\frac{w}{1+w},\qquad
 L=3+t\E_{\nu_t}X,\qquad
 V=t^2\operatorname{Var}_{\nu_t}(X).
\]
The mean under the tilt is increasing in $t$ and starts at
$(\lambda_1+\lambda_2+\lambda_3)/4=(a+b)/4\geq0$; hence $L\geq3$.
Using $t\,d/dt$ in \eqref{eq:real-first-simplex} gives the exact identity
\begin{equation}\label{eq:real-first-simplex-curvature}
 -t^2(\log h_1)''
 =\theta\bigl(3-V-(1-\theta)L^2\bigr).
\end{equation}
Consequently the first-start target is already proved whenever
\[
 \theta\leq\frac{1+\sqrt3}{3},
 \qquad\text{equivalently}\qquad
 w\leq5+3\sqrt3,
\]
because $V\geq0$ and $L\geq3$ make the right side of
\eqref{eq:real-first-simplex-curvature} at most
$\theta(9\theta-6)\leq2$.  In the complementary range, the whole real-root
problem is the scalar tilted-simplex inequality
\begin{equation}\label{eq:real-first-simplex-target}
 V+(1-\theta)L^2\geq3-\frac2\theta.
\end{equation}
Thus no signed partial-fraction coefficients remain in the real-root target.

The real-root region also has a compact polynomial chart, and on a substantial
initial range the target admits an exact finite certificate.  After scaling,
write the three roots as
\[
 m-\rho,\qquad m+\rho,\qquad m+1.
\]
The coefficient identities for $P$ and the nonnegativity of $a,b,Q$ show that
all three roots are nonnegative and that the admissible set is exactly
\begin{equation}\label{eq:real-root-compact-chart}
 0\leq m\leq1,\qquad
 \max\left\{0,\frac{3m^2+2m-1}{4}\right\}
 \leq\rho^2\leq m^2.
\end{equation}
Indeed, if the two smaller roots were negative, their sum with the Perron
root and their second elementary symmetric function could not both be
nonnegative.  The discriminant of the quadratic with roots $a,b$ gives the
lower bound in \eqref{eq:real-root-compact-chart}.  This region is the union
of the two polynomial images of the unit square
\begin{align}
 m&=\frac r3,&
 \rho^2&=\tau m^2,\label{eq:real-root-low-rectangle}\\
 m&=\frac{1+2r}{3},&
 \rho^2&=\frac{r(r+2)}3+\tau\frac{(1-r)^2}{9},
 \label{eq:real-root-high-rectangle}
\end{align}
where $0\leq r,\tau\leq1$.

Let $a_n=h_1^{(n)}(0)$ in this normalization.  Then
\[
 a_0=1,\qquad a_1=a_2=0,
\]
and
\begin{equation}\label{eq:real-first-derivative-recurrence}
 a_{n+3}=(1+3m)a_{n+2}-(3m^2+2m-\rho^2)a_{n+1}
 +(1+m)(m^2-\rho^2)a_n.
\end{equation}
If $D_N$ denotes $N!$ times the coefficient of degree $N$ in the corrected
curvature, then
\begin{equation}\label{eq:real-first-taylor-coefficient}
 D_N=\sum_{j=0}^N\binom Nj
 \left(2-\frac N2+\frac{(N-2j)^2}{2}\right)a_ja_{N-j}.
\end{equation}
Exact conversion shows that every $D_N$, $0\leq N\leq60$, has only
nonnegative tensor Bernstein coefficients on each of
\eqref{eq:real-root-low-rectangle} and
\eqref{eq:real-root-high-rectangle}.

This finite calculation already gives a rigorous interval, not merely
coefficient evidence.  Normalize the Perron root to $1$ and call the other
two roots $0\leq u\leq v\leq1$.  Positive rescaling of time multiplies
$D_N$ by a positive factor and therefore preserves the Bernstein-certified
signs.  In the new normalization,
\[
 \sqrt u+\sqrt v\leq1,
 \qquad
 a_n=uv\,h_{n-3}(u,v,1)\quad(n\geq3),
\]
where $h_k$ is the complete homogeneous polynomial.  Consequently
\[
 0\leq a_n\leq\frac{uv}{(1-u)(1-v)}\leq\frac19.
\]
For $N\geq61$, the absolute value of
\eqref{eq:real-first-taylor-coefficient} is therefore at most $N^2 2^N$.
On $0\leq x\leq10$, where $x$ now denotes Perron-normalized time, the tail
after degree $60$ is at most
\[
 \frac{61^2 20^{61}}{61!}
 \left(1-\frac{20\cdot62}{61^2}\right)^{-1}<1.
\]
The Bernstein partial sum is at least its constant term $2$, so the
first-start inequality is proved throughout this compact time range.  The
calculation and the rational tail comparison are reproduced by
\path{python/check_real_first_full.py}.

For reference, the unbounded tail has a particularly short exact form.
For $0<u<v<1$, partial fractions give
\[
 h_1(t)=\frac{v e^{ut}}{(v-u)(1-u)}
 -\frac{u e^{vt}}{(v-u)(1-v)}
 +\frac{uv e^t}{(1-u)(1-v)}.
\]
Put
\[
 A=(1-v)t,\qquad B=(1-u)t,
\]
and
\[
 R=\frac{1-u}{v(v-u)}e^{-A},\qquad
 S=\frac{1-v}{u(v-u)}e^{-B}.
\]
After factoring the last exponential, direct differentiation yields
\begin{equation}\label{eq:real-first-tail-energy}
 -t^2(\log h_1)''
 =\frac{RA^2+RS(B-A)^2-SB^2}{(1-R+S)^2}.
\end{equation}
The apparent positive mixed term is already controlled by the last term:
\begin{equation}\label{eq:real-first-tail-mixed-control}
 R(B-A)^2\leq B^2.
\end{equation}
Indeed, after cancelling $t^2$, this follows from
$v-u\leq v(1-u)$ and $e^{-A}\leq1$.  Thus the only part not covered by
the compact certificate is proving that the right side of
\eqref{eq:real-first-tail-energy} is at most $2$ for $t\geq10$, including
the confluent limits.  This is now a two-exponential tail estimate on the
compact root simplex, rather than a matrix inequality.  The partial-fraction
identity, \eqref{eq:real-first-tail-energy}, and
\eqref{eq:real-first-tail-mixed-control} are checked symbolically by
\path{python/check_real_first_tail_identities.py}.

There is a sharper coefficient route to the tail.  Put
$g_n=h_n(u,v,1)$ and consider the normalized complete homogeneous sequence
\begin{equation}\label{eq:real-first-normalized-complete-target}
 q_n=\frac{g_n}{n+1}.
\end{equation}
The remaining sequence statement is elementary, although its useful form is
not the most familiar normalization.
\begin{proposition}[Normalized complete-homogeneous log-convexity]
\label{prop:real-first-normalized-complete-logconvex}
\begin{equation}\label{eq:real-first-normalized-complete-logconvex}
 q_n^2\leq q_{n-1}q_{n+1}\qquad(n\geq1)
 \quad\hbox{whenever}\quad \sqrt u+\sqrt v\leq1.
\end{equation}
\end{proposition}
Notice that the case $n=1$ is precisely
$(1-u-v)^2\geq4uv$.
The Jacobi--Trudi identity gives an equivalent symmetric-polynomial form:
\begin{equation}\label{eq:real-first-schur-target}
 n(n+2)h_n(u,v,uv)
 \leq h_{n-1}(1,u,v)h_{n+1}(1,u,v).
\end{equation}
Indeed,
$s_{(n,n)}(1,u,v)=g_n^2-g_{n-1}g_{n+1}=h_n(u,v,uv)$.
This formulation may be more convenient for a uniform proof because it
contains no quotients and its $n=1$ case is exactly the root discriminant.

There is now stronger exact evidence for a coefficientwise proof.  Introduce
nonnegative $x,y,r$ by
\[
 u=x^2,\qquad v=y^2,\qquad 1=(x+y)^2+r
\]
after homogeneous rescaling, and put
\[
 \mathcal D_n(A,B,C)
 =(n+1)^2h_{n-1}(A,B,C)h_{n+1}(A,B,C)
   -n(n+2)h_n(A,B,C)^2.
\]
Exact coefficient extraction gives, through $n=40$,
\begin{equation}\label{eq:real-first-coefficient-cone}
 \mathcal D_n\bigl(x^2,y^2,(x+y)^2+r\bigr)
 \in \mathbb Z_{\geq0}[x,y,r].
\end{equation}
It also gives the apparently stronger monotonicity
\begin{equation}\label{eq:real-first-coefficient-increment}
 \mathcal D_{n+1}\bigl(x^2,y^2,C\bigr)
 -C^2\mathcal D_n\bigl(x^2,y^2,C\bigr)
 \in \mathbb Z_{\geq0}[x,y,r],
 \qquad C=(x+y)^2+r.
\end{equation}
For example,
\[
 \mathcal D_1\bigl(x^2,y^2,(x+y)^2+r\bigr)=r(r+4xy).
\]
Thus a uniform proof of either coefficient assertion would prove
\eqref{eq:real-first-normalized-complete-logconvex} immediately, and
\eqref{eq:real-first-coefficient-increment} suggests a particularly short
induction.  The integer coefficient calculation avoids symbolic expansion and
is reproduced by \path{python/check_normalized_complete_cone.py}.

We now prove the proposition.  By continuity, order the roots as
$0<u<v<1$; the confluent and boundary cases follow by taking limits.  Let
$X$ be uniform on $[v,1]$
and $Y$ uniform on $[u,v]$, and write $A_j=\mathbb E X^j$ and
$B_j=\mathbb E Y^j$.  Partial fractions followed by one integration gives
\begin{equation}\label{eq:real-first-uniform-difference}
 q_n=\frac{A_n-uB_n}{1-u}.
\end{equation}
Indeed, the representing signed density vanishes on $[0,u]$, equals
$-u(1-v)/(v-u)$ on $[u,v]$, and equals $1$ on $[v,1]$, before division by
$(1-u)(1-v)$.  Consequently
\begin{align}
 (1-u)^2(q_{n-1}q_{n+1}-q_n^2)
  &=\Delta_A+u^2\Delta_B-u\Gamma,\label{eq:real-first-uniform-determinant}\\
 \Delta_A&=A_{n-1}A_{n+1}-A_n^2,\qquad
 \Delta_B=B_{n-1}B_{n+1}-B_n^2,\nonumber\\
 \Gamma&=\mathbb E\!\left[X^{n-1}Y^{n-1}(X-Y)^2\right].\nonumber
\end{align}
Here both $\Delta$ terms are nonnegative.  The case $n=1$ is the root
discriminant noted above.  For $n\geq2$, pointwise
maximization in $Y$ gives
\begin{equation}\label{eq:real-first-cross-bound}
 \Gamma\leq
 \min\!\left\{v^{n-1},
 \frac{4(n-1)^{n-1}}{(n+1)^{n+1}}\right\}A_{n+1}.
\end{equation}
Since the root condition implies
$u\leq\min\{v,(1-\sqrt v)^2\}$, it remains only to prove the
one-variable inequality
\begin{equation}\label{eq:real-first-one-variable-moment-bound}
 \min\{v,(1-\sqrt v)^2\}
 \min\!\left\{v^{n-1},
 \frac{4(n-1)^{n-1}}{(n+1)^{n+1}}\right\}A_{n+1}
 \leq\Delta_A .
\end{equation}
Put
\[
 c_n=\frac{4(n-1)^{n-1}}{(n+1)^{n+1}},
 \qquad
 A_j=\frac{1-v^{j+1}}{(j+1)(1-v)}.
\]
Direct simplification gives
\begin{equation}\label{eq:uniform-moment-determinant-polynomial}
 \Delta_A
 =\frac{(1-v)^2P_n(v)}{n(n+1)^2(n+2)},
\end{equation}
where
\begin{equation}\label{eq:uniform-moment-palindromic-polynomial}
 P_n(v)=
 \sum_{k=0}^{n-2}\binom{k+3}{3}
       \bigl(v^k+v^{2n-2-k}\bigr)
 +\binom{n+2}{3}v^{n-1}.
\end{equation}
Thus $P_n(v)\geq1$.  Pairing the terms in each parenthesis by AM--GM and
using $\sum_{k=0}^{n-2}\binom{k+3}{3}=\binom{n+2}{4}$ also gives
\begin{equation}\label{eq:uniform-moment-central-lower-bound}
 P_n(v)\geq\frac{n(n+1)^2(n+2)}{12}\,v^{n-1}.
\end{equation}

First suppose $0\leq v\leq1/4$ and $n\geq6$.  Since
$A_{n+1}\leq((n+2)(1-v))^{-1}$, the left side of
\eqref{eq:real-first-one-variable-moment-bound} is at most
$v^n/((n+2)(1-v))$.  Equations
\eqref{eq:uniform-moment-determinant-polynomial} and $P_n\geq1$ reduce the
claim to
\[
 n(n+1)^2v^n\leq(1-v)^3.
\]
It is enough to set $v=1/4$.  The resulting inequality holds already at
$n=4$, and its left side decreases with $n$.

Now suppose $1/4\leq v\leq1$, write $b=\sqrt v$, and take $n\geq6$.
By \eqref{eq:uniform-moment-central-lower-bound}, it is enough to show
\begin{equation}\label{eq:uniform-moment-transition-bound}
 \frac{12A_{n+1}}{(1+b)^2}
 \min\left\{1,\frac{c_n}{b^{2n-2}}\right\}\leq1.
\end{equation}
Set
\[
 b_n=c_n^{1/(2n-2)}
 =\sqrt{\frac{n-1}{n+1}}
   \left(\frac2{n+1}\right)^{1/(n-1)}.
\]
For $1/2\leq b\leq b_n$, use
$A_{n+1}\leq((n+2)(1-b^2))^{-1}$.  Since
$(1-b)(1+b)^3$ decreases for $b\geq1/2$, it suffices that
\begin{equation}\label{eq:uniform-moment-bn-bound}
 (n+2)(1-b_n)(1+b_n)^3\geq12.
\end{equation}
For $b\geq b_n$, the left side of
\eqref{eq:uniform-moment-transition-bound} is nonincreasing.  Indeed,
\[
 A_{n+1}=\frac1{n+2}\sum_{j=0}^{n+1}b^{2j},
\]
and $bA_{n+1}'/A_{n+1}\leq n+1$, because the weights $b^{2j}$ decrease
with $j$.  This is at most
$2n-2+2b/(1+b)$, the logarithmic derivative with respect to $\log b$ of
$b^{2n-2}(1+b)^2$.  Thus this region also reduces to
\eqref{eq:uniform-moment-bn-bound}.

For completeness, \eqref{eq:uniform-moment-bn-bound} has a short calculus
check.  Extend $b_n$ to $b(x)$ for real $x\geq6$.  Then
\[
 \frac{d}{dx}\log b(x)
 =\frac{\log((x+1)/2)}{(x-1)^2}>0.
\]
Moreover, $(x+2)(1-b(x))$ is increasing.  To see this, use
$b^{-1}-1\geq-\log b$ and
\[
 \frac{(x-1)^2}{2}\log\frac{x+1}{x-1}
 \geq x-2+\frac1x
 \geq3\log\frac{x+1}{2}\qquad(x\geq6).
\]
Hence $(x+2)(1-b(x))(1+b(x))^3$ is increasing.  Finally
$b(6)<2/3$, so its value at $x=6$ is greater than
$8(1-2/3)(1+2/3)^3=1000/81>12$.

It remains only to cover $n=2,3,4,5$.  On the two intervals
$0\leq b\leq1/2$ and $1/2\leq b\leq1$, exact Bernstein conversion proves
the stronger inequality
\[
 \Delta_A\geq \min\{b^2,(1-b)^2\}\,c_nA_{n+1}.
\]
The rational certificates are reproduced by
\path{python/check_uniform_moment_bound.py}.  This proves
\eqref{eq:real-first-one-variable-moment-bound}; inserting it into
\eqref{eq:real-first-uniform-determinant} proves the proposition.

Here is why the proposition finishes the tail.  Write
$w=h_1-1$ and
\[
 t^2\{ww''-(w')^2\}+2w^2
 =(uv)^2\sum_{N\geq6}\widehat D_N\frac{t^N}{N!}.
\]
Since $w=uv\sum_{N\geq3}g_{N-3}t^N/N!$, one has, with
$\kappa_N(j)=2-N/2+(N-2j)^2/2$,
\[
 \widehat D_N=
 \sum_{j=3}^{N-3}\binom Nj\kappa_N(j)
 g_{j-3}g_{N-j-3}.
\]
Set $M=N-6$ and
\[
 B_i=\binom N{i+3}(i+1)(M-i+1),\qquad 0\leq i\leq M.
\]
The summand weight is $B_iq_iq_{M-i}$.  By
\eqref{eq:real-first-normalized-complete-logconvex}, the symmetric
multiplier $q_iq_{M-i}$ is nondecreasing with
$|i-M/2|$.  The same is true of $\kappa_N(i+3)$.  The weighted Chebyshev
inequality therefore compares its average with the one obtained by setting
$q_i=1$.  The numerator of that baseline average is
\begin{equation}\label{eq:real-first-baseline-coefficient}
 \sum_{i=0}^M B_i\kappa_N(i+3)
 =\frac{2^N(N^2-17N+32)+4N^4-24N^3+68N^2-48N-64}{4},
\end{equation}
which is positive for every $N\geq15$.  Hence the proposition implies
$\widehat D_N\geq0$ for all $N\geq15$.

The finitely many lower coefficients already combine positively in the
tail range.  With $y=10/t$, exact tensor Bernstein conversion on
$0\leq s,d,y\leq1$, where
\[
 \sqrt u=\frac{s(1-d)}2,\qquad
 \sqrt v=\frac{s(1+d)}2,
\]
proves
\begin{equation}\label{eq:real-first-variance-low-block}
 t^{-16}\sum_{N=6}^{16}\widehat D_N\frac{t^N}{N!}\geq0
 \qquad(t\geq10).
\end{equation}
Thus \eqref{eq:real-first-normalized-complete-logconvex} implies
$V\geq1$ on the whole remaining tail, and
\eqref{eq:real-first-simplex-curvature} gives
$E\leq\theta(2-(1-\theta)L^2)\leq2$.  The exact low block, the signs of
$\widehat D_N$ through $N=40$, and the first twenty cases of the sequence
lemma are reproduced by
\path{python/check_real_first_variance_tail.py}.

For comparison, a second possible reduction is suggested by scale
monotonicity.  Put
\[
 \sqrt u=\frac{s(1-d)}2,\qquad
 \sqrt v=\frac{s(1+d)}2,\qquad 0\leq s,d\leq1.
\]
Numerical differentiation indicates that $E$ is nondecreasing in $s$
whenever $E\geq0$.  Such a conditional monotonicity would reduce the tail
to $s=1$, which is exactly the repeated-rate boundary $a=b$.  The observed
maximum on that boundary is approximately $1.30148$.  These observations
are reproduced by \path{python/explore_real_first.py}; they are evidence,
not part of the proof.

The complex-pair region itself admits a compact two-parameter normalization.
This removes the three original rates from a direct analytic attack on the
first start.  Write the roots of $P$ as
$\lambda,\mu\mathbin\pm i\nu$, put $d=\lambda-\mu>0$, and set
\[
 x=dt,\qquad m=\frac{\mu}{d},\qquad \omega=\frac{\nu}{d}.
\]
Comparison of the coefficients of $P$ gives
\[
 a+b=d(1+3m),\qquad
 ab=d^2(3m^2+2m+\omega^2),\qquad
 Q=d^3(1+m)(m^2+\omega^2).
\]
Consequently the admissible parameter set is
\begin{equation}\label{eq:complex-root-domain}
 -\frac13\leq m\leq\frac13,
 \qquad
 \max\{0,-3m^2-2m\}\leq\omega^2
 \leq\frac{1-2m-3m^2}{4},
\end{equation}
with degenerate boundary points understood by continuity.  After the harmless
exponential shift $e^{-\mu t}$, the first component has the explicit form
\begin{equation}\label{eq:complex-root-first-start}
 e^{-\mu t}h_1(t)=\frac{H_{m,\omega}(x)}{1+\omega^2},
\end{equation}
where
\[
 H_{m,\omega}(x)
 =(m^2+\omega^2)e^x+(1-m^2)\cos(\omega x)
 -\frac{(m+1)(m+\omega^2)}{\omega}\sin(\omega x).
\]
Thus the first bound in the complex-pair region is exactly the elementary
three-variable inequality
\begin{equation}\label{eq:complex-root-curvature}
 x^2\bigl(H_{m,\omega}H_{m,\omega}''-(H_{m,\omega}')^2\bigr)
 +2H_{m,\omega}^2\geq0
\end{equation}
for $x\geq0$ and the compact parameter domain
\eqref{eq:complex-root-domain}.  This is an equivalent reduction, not yet a
proof, but it supplies a second concrete target with no matrix or
Laplace-transform variables left.

The Wronskian in \eqref{eq:complex-root-curvature} admits a further exact
factorization.  This removes the second derivatives and isolates the only
term which can be negative.  Put
\[
 z=\omega^2,\qquad A=m^2+z,\qquad
 T=(1-m^2)\cos(\omega x)-(m+1)(m+z)S_\omega(x),
\]
so that $H_{m,\omega}=Ae^x+T$.  The trigonometric part is a harmonic
oscillator, and direct simplification gives
\begin{align}
 T'^2+zT^2&=A(m+1)^2(1+z),\label{eq:first-oscillator-energy}\\
 (1-z)T-2T'
 &=(m+1)(1+z)
 \left((m+1)\cos(\omega x)+(z-m)S_\omega(x)\right).
 \label{eq:first-oscillator-companion}
\end{align}
Consequently
\begin{align}
 H_{m,\omega}H_{m,\omega}''
 -(H_{m,\omega}')^2
 ={}&A(m+1)(1+z)\bigl[
 e^x\{(m+1)\cos(\omega x)+(z-m)S_\omega(x)\}\notag\\
 &\hspace{43mm}{}-(m+1)\bigr].
 \label{eq:first-oscillator-wronskian}
\end{align}
Thus the complex first-start target is the derivative-free inequality
\begin{align}
 2(Ae^x+T)^2
 +x^2A(m+1)(1+z)\bigl[
 e^x\{(m+1)\cos(\omega x)+(z-m)S_\omega(x)\}
 -(m+1)\bigr]\geq0.
 \label{eq:first-oscillator-target}
\end{align}
All factors outside the square brackets are nonnegative on
\eqref{eq:complex-root-domain}.  At the confluent boundary $z=0$ this has
the continuous specialization
\[
 H_{m,0}(x)=m^2e^x+1-m^2-m(m+1)x,\qquad 0\leq m\leq\frac13,
\]
and the bracket in \eqref{eq:first-oscillator-target} becomes
$e^x(m+1-mx)-(m+1)$.  These identities are checked symbolically by
\texttt{python/check\_complex\_first\_factorization.py}.

The confluent boundary just displayed can in fact be settled completely.
This is one of the two lower boundary arcs suggested by
\eqref{eq:first-oscillator-target}.
\begin{proposition}[Confluent complex boundary]\label{prop:first-omega-zero}
For $z=0$, $0\leq m\leq1/3$, and $x\geq0$, the first-start inequality
\eqref{eq:complex-root-curvature} holds.
\end{proposition}
\begin{proof}
The case $m=0$ is immediate, so suppose $m>0$ and write $r=1+m$.  With
\[
 H=m^2e^x+1-m^2-mrx,\qquad
 J=e^x(r-mx)-r,
\]
the left side of \eqref{eq:complex-root-curvature} is
\begin{equation}\label{eq:first-omega-zero-D}
 D=2H^2+x^2m^2rJ.
\end{equation}
If $J\geq0$ there is nothing to prove.  Since
$J'=e^x(1-mx)$ and $J(0)=0$, the alternative $J<0$ forces $x>1/m$.

We use twice the elementary bound
\[
 \frac{m^2e^{1/m}}{1+m}\geq\frac{e^3}{12}>\sqrt2
 \qquad(0<m\leq1/3).
\]
Indeed, the expression on the left is decreasing on this interval.  It
first implies, for $x\geq1/m$,
\[
 H\geq\frac12m^2e^x.
\]
If moreover $x\leq r/m$, then $-J\leq r$ and
$me^x\geq\sqrt2\,xr$; hence
\[
 D\geq\frac12m^4e^{2x}-x^2m^2r^2\geq0.
\]

It remains to consider $x\geq r/m$.  Regard the right side of
\eqref{eq:first-omega-zero-D} as the quadratic
\[
 \mathcal D(m,x,Y)
 =2(m^2Y+1-m^2-mrx)^2+x^2m^2r\{Y(r-mx)-r\}.
\]
Put $u=mx-r\geq0$.  Then
\[
 m=\frac{u+1}{x-1},\qquad x=3u+4+v
\]
for some $v\geq0$, because $m\leq1/3$.  Also $x\geq4$ and
$e^x\geq x^4/6$: the quotient $e^x/x^4$ is increasing there, and the
inequality at $x=4$ follows from the first six terms of the exponential
series.  Direct expansion after the preceding substitution shows that both
\[
 72(x-1)^4\mathcal D\left(m,x,\frac{x^4}{6}\right),
 \qquad
 6(x-1)^4\partial_Y\mathcal D\left(m,x,\frac{x^4}{6}\right)
\]
are polynomials in $u,v$ with nonnegative integer coefficients.  Therefore
$\mathcal D$ is nonnegative at $Y=x^4/6$ and is increasing for all larger
$Y$.  Taking $Y=e^x$ proves the result.  The two polynomial expansions
(respectively $80$ and $35$ monomials) are verified exactly by
\path{python/check_complex_first_factorization.py}.
\end{proof}

The other lower boundary arc can also be proved.  The following argument uses
an exact finite Bernstein certificate only on a compact interval; the tail is
handled analytically.
\begin{proposition}[$a=0$ complex boundary]
\label{prop:first-a-zero}
For $m=-q$, $0\leq q\leq1/3$, and
$z=2q-3q^2$, the first-start inequality
\eqref{eq:complex-root-curvature} holds for every $x\geq0$.
\end{proposition}
\begin{proof}
Undo the harmless exponential shift and put
\[
 h(x)=e^{-qx}\frac{H_{-q,\sqrt z}(x)}{1+z}.
\]
In the ordered canonical coordinates this is the first component with
\[
 b=1-3q,\qquad Q=2q(1-q)^2,
\]
and hence
\begin{equation}\label{eq:a-zero-scalar-ode}
 h'''=(1-3q)h''+2q(1-q)^2h,
 \qquad h(0)=1,\quad h'(0)=h''(0)=0.
\end{equation}
Multiplication by $e^{-qx}$ does not change logarithmic curvature, so it is
enough to prove
\[
 D_q(x)=x^2\{hh''-(h')^2\}+2h^2\geq0.
\]

We first treat $0\leq x\leq10$.  Write
$a_n(q)=h^{(n)}(0)$.  Equation~\eqref{eq:a-zero-scalar-ode} gives the exact
polynomial recurrence
\[
 a_{n+3}=(1-3q)a_{n+2}+2q(1-q)^2a_n,
 \qquad a_0=1,\quad a_1=a_2=0.
\]
If $D_N(q)$ is $N!$ times the coefficient of $x^N$ in $D_q(x)$, then
\[
 D_N(q)=\sum_{j=0}^N\binom Nj
 \left(2-\frac N2+\frac{(N-2j)^2}{2}\right)
 a_j(q)a_{N-j}(q).
\]
Exact rational calculation shows that, for every $0\leq N\leq80$,
$D_N(x/3)$ has nonnegative coefficients in the Bernstein basis on
$0\leq x\leq1$.  Thus the degree-$80$ Taylor polynomial of $D_q$ is at
least its constant term $2$ throughout the required parameter interval.

For completeness, the Taylor tail has a short uniform bound.  The spectral
form is
\begin{equation}\label{eq:a-zero-spectral-form}
 h=c_0e^{(1-q)x}+c_+e^{(-q+i\sqrt z)x}
   +c_-e^{(-q-i\sqrt z)x},
 \qquad c_-=\overline{c_+},
\end{equation}
where
\[
 c_0=\frac{2q(1-q)}{1+z}.
\]
All three exponents have modulus at most $1$, and
$|c_0|+|c_+|+|c_-|\leq2$.  For a sum
$h=\sum c_re^{\lambda_rx}$ one has
\[
 D_q(x)=\sum_{r,s}c_rc_s
 \left\{2+\frac{x^2}{2}(\lambda_r-\lambda_s)^2\right\}
 e^{(\lambda_r+\lambda_s)x}.
\]
Consequently, on $0\leq x\leq10$, the absolute remainder after total degree
$80$ is at most
\[
 8e^{20}\left\{\frac{20^{81}}{81!}
 +100\frac{20^{79}}{79!}\right\}<1.
\]
Here one may use the elementary bound $e^{20}<500000000$; the final
inequality is rational arithmetic.  The exact Bernstein calculation and
this tail check are reproduced by
\path{python/check_complex_first_a0_boundary.py}.  It follows that
$D_q(x)>1$ on this compact interval.

It remains to give a purely analytic estimate for $x\geq10$.  Put
\[
 U=\cos(\sqrt z,x)+3qS_{\sqrt z}(x).
\]
On the present boundary the square bracket in
\eqref{eq:first-oscillator-target} is $(1-q)(e^xU-1)$.  If it is
nonnegative, the result is immediate.  Otherwise $U<e^{-x}<1/2$.  If
$\sqrt z,x\leq1$, then
$\cos(\sqrt z,x)\geq1-zx^2/2\geq1/2$ and the sine term in $U$ is
nonnegative, a contradiction.  Hence $zx^2>1$.  Since $z\leq2q$, this
implies
\begin{equation}\label{eq:a-zero-q-lower}
 q>\frac1{2x^2},\qquad c_0\geq q>\frac1{2x^2}.
\end{equation}

Let $R=|c_+|+|c_-|$.  From the cosine and sine coefficients in
\eqref{eq:a-zero-spectral-form},
\[
 R\leq\frac{1-q^2}{1+z}
 +\frac{q(1-q)(1-3q)}{\sqrt z(1+z)}\leq\frac54.
\]
For the last inequality, the first summand is at most $1$, while the square
of the assertion that the second is at most $1/4$ reduces to
\[
 171q^3-96q^2+7q-2\leq0.
\]
This follows from $171q^3\leq57q^2$ and the negative-discriminant quadratic
$-39q^2+7q-2$.

The leading--oscillatory spectral gaps have squared modulus at most $4/3$.
Therefore, with $L=2+2x^2/3$, formula
\eqref{eq:a-zero-spectral-form} gives
\[
 D_q(x)\geq e^{-2qx}\bigl[
 2c_0^2e^{2x}-L\{2c_0Re^x+R^2\}\bigr].
\]
For $x\geq10$ the bracket is increasing as a function of
$c_0\geq1/(2x^2)$: its derivative is nonnegative because
$e^x/x^2\geq LR$, which follows already from
$L\leq7x^2/10$, $R\leq5/4$, and $e^{10}>22000$.
Using the same bounds, its value at the lower endpoint is at least
\[
 \frac{e^{2x}}{2x^4}-\frac78e^x-\frac{35}{32}x^2>0.
\]
Indeed, after division by $e^x$, the first term is increasing and the last
is decreasing for $x\geq10$; at $x=10$ the claim follows already from
$e^{10}>22000$.  This proves the tail and completes the boundary.
\end{proof}

In fact the same finite-certificate strategy, with a uniform complex-analytic
remainder estimate, fills the interior of the complex-root chart.
\begin{proposition}[Complex-pair first start]
\label{prop:complex-first}
For every parameter pair in \eqref{eq:complex-root-domain} and every
$x\geq0$, the first-start inequality
\eqref{eq:complex-root-curvature} holds.
\end{proposition}
\begin{proof}
We first prove the assertion on $0\leq x\leq10$.  Two overlapping choices of
the harmless exponential shift keep the exact certificates small.  For
$m\geq0$, use
\[
 F=\frac{H_{m,\omega}}{1+z}.
\]
Its derivatives at zero satisfy
\[
 f_0=1,\qquad f_1=-m,\qquad f_2=m^2,
 \qquad f_{n+3}=f_{n+2}-zf_{n+1}+zf_n.
\]
Let $P_{70}$ be the degree-$70$ Taylor polynomial of
$x^2\{FF''-(F')^2\}+2F^2$.  Parameterize the positive half of
\eqref{eq:complex-root-domain} by
\[
 m=\frac r3,\qquad
 z=\tau\frac{3-2r-r^2}{12},\qquad 0\leq r,\tau\leq1,
\]
and put $x=10u$, $0\leq u\leq1$.  Exact conversion of
$P_{70}(10u)$ to the tensor Bernstein basis in $(r,\tau,u)$ gives only
positive coefficients; the smallest is
\begin{equation}\label{eq:complex-first-bernstein-min}
 \frac{7445875529}{6372236871}>1.
\end{equation}

For $m\leq0$, use instead
$h=e^{mx}F$, the unshifted canonical component.  Its derivatives satisfy
\[
 h_{n+3}=(1+3m)h_{n+2}-(3m^2+2m+z)h_{n+1}
 +(1+m)(m^2+z)h_n,
 \qquad h_0=1,\qquad h_1=h_2=0.
\]
Write $r=-3m$ and parameterize this half by
\[
 z=\frac{r(2-r)}3+\tau\frac{(1-r)^2}{4},\qquad
 0\leq r,\tau\leq1.
\]
For every Taylor degree $0\leq N\leq70$, the corresponding curvature
coefficient has only nonnegative tensor Bernstein coefficients in
$(r,\tau)$.  Hence the degree-$70$ partial sum is at least its constant
term $2$ for $x\geq0$.

The omitted tail is uniformly smaller than $1$ in both cases.  Indeed, the
chosen function ($F$ for $m\geq0$ and $h$ for $m\leq0$), together with its
first two derivatives, has modulus at most $36e^{34}$ on the circle
$|x|=34$.  This follows directly from the exponential--sine--cosine formula:
all relevant spectral frequencies have modulus at most $1$, all scalar
coefficients have modulus at most $1$, and
$|S_\omega(x)|\leq |x|e^{|\omega x|}$.  Thus the curvature has modulus at
most
\[
 M=(2\cdot34^2+2)(36e^{34})^2.
\]
Cauchy's estimate bounds its Taylor tail on $|x|\leq10$ by
\[
 M\frac{(5/17)^{71}}{1-5/17}<1;
\]
the last inequality follows already from $e^{34}<6\cdot10^{14}$ and is an
exact rational check.  The two Bernstein certificates and the tail bound
are verified by \path{python/check_complex_first_full.py}.  Combined with
\eqref{eq:complex-first-bernstein-min}, they prove the claim for
$0\leq x\leq10$.

It remains to give a uniform analytic argument for $x\geq10$.  Put
\[
 A=m^2+z,\qquad
 U=\cos(\omega x)+\frac{z-m}{m+1}S_\omega(x),\qquad y=x\sqrt A.
\]
The square bracket in \eqref{eq:first-oscillator-target} is
$(m+1)(e^xU-1)$.  If it is nonnegative, there is nothing to prove.  Suppose
therefore that $U<e^{-x}$.  The parameter domain gives
\[
 z\leq\frac13,\qquad A\leq\frac49,\qquad
 (z-m)^2\leq A,\qquad m+1\geq\frac23.
\]
Using $\cos v\geq1-v^2/2$ and $|S_\omega(x)|\leq x$, we obtain
\[
 U\geq1-\frac{y^2}{2}-\frac{3y}{2}.
\]
Since $e^{-x}<1/8$, this forces $y>1/2$.

The trigonometric term in $H=Ae^x+T$ has the useful identity
\begin{equation}\label{eq:complex-first-TU}
 T=(1-m^2)U-2A S_\omega(x).
\end{equation}
If $0\leq U<e^{-x}$, then
$H\geq A(e^x-2x)\geq(9/10)Ae^x$.  If $U<0$, put
\[
 G(y)=\frac{y^2}{2}+\frac{3y}{2}-1.
\]
The preceding lower bound gives $|U|\leq G(y)$, while
\[
 G(y)\leq\frac{17}{16}y^2,
\]
because the difference after division by $y^2$ is
$(3y-4)^2/(16y^2)$.  Equation~\eqref{eq:complex-first-TU} then gives
\[
 H\geq Ae^x-G(y)-2Ax\geq\frac9{10}Ae^x;
\]
the last step uses $x\geq10$, $y>1/2$, and $e^{10}>22000$.

Finally,
$(m+1)^2(1+z)\leq64/27$.  In the case $U\geq0$, the possibly negative
bracket has magnitude at most $1$, and hence
\[
 D_H\geq\frac32A^2e^{2x}-\frac{64}{27}x^2A>0.
\]
In the case $U<0$, use
$1+e^x|U|\leq1+(17/16)e^xy^2$ to obtain
\[
 D_H\geq y^4\left\{
 \frac{3e^{2x}}{2x^4}-\frac{68}{27}e^x-\frac{256}{27}
 \right\}>0.
\]
Both final inequalities follow at $x=10$ from $y>1/2$ and
$e^{10}>22000$, and their left sides only improve thereafter.  This proves
the full complex-pair region.
\end{proof}

There is also a natural, but no longer necessary, monotonicity route from the
two boundaries into the interior of the
complex-root domain.  Normalize
$\widehat H=H_{m,\omega}/(1+z)$.  Its Laplace transform with respect to $x$
and its $z$-derivative are
\begin{align}
 \mathcal L\widehat H(s)
 &=\frac{s^2-(m+1)s+m^2+m+z}{(s-1)(s^2+z)},\notag\\
 \partial_z\mathcal L\widehat H(s)
 &=\frac{(m+1)(s-m)}{(s-1)(s^2+z)^2}.
 \label{eq:first-z-derivative-laplace}
\end{align}
Equivalently,
\[
 \partial_z\widehat H
 =(m+1)\,e^x*
 \bigl(\cos(\omega x)-mS_\omega(x)\bigr)*S_\omega(x).
\]
Exact numerical minimization indicates that the normalized curvature in
\eqref{eq:complex-root-curvature}, divided by $(1+z)^2$, is nondecreasing in
$z$ for fixed $(x,m)$.  This monotonicity is no longer needed for the proof,
and its individual Taylor coefficients do change sign.  Formula
\eqref{eq:first-z-derivative-laplace} is a compact starting point for that
stronger comparison; at present it remains conjectural.

The missing middle bound has an equally explicit form in the same chart.
Let
\[
 \delta=\sqrt{1-2m-3m^2-4\omega^2},\qquad
 \beta=\frac{1+m+\delta}{2}=\frac{b-\mu}{d}.
\]
Then, up to the positive constant $x_2/d$,
\[
 e^{-\mu t}h_2(t)=\frac{K_{m,\omega}(x)}{1+\omega^2},
\]
where
\[
 K_{m,\omega}(x)
 =(1-\beta)\bigl(e^x-\cos(\omega x)\bigr)
 +\frac{\omega^2+\beta}{\omega}\sin(\omega x).
\]
Thus the sharper middle target in \eqref{eq:two-riccati-targets} is equivalent
in the complex-pair region to
\begin{equation}\label{eq:complex-root-middle-curvature}
 x^2\bigl(K_{m,\omega}K_{m,\omega}''-(K_{m,\omega}')^2\bigr)
 +K_{m,\omega}^2\geq0.
\end{equation}
The quotients by $\omega$ in both formulas have their continuous meanings at
$\omega=0$.

The parameter dependence of the last inequality can be reduced further.  Set
\[
 A=1-\beta=\frac{\lambda-b}{d},\qquad
 \rho=\sqrt{1-3\omega^2}.
\]
Eliminating $m$ and $b$ from \eqref{eq:complex-root-domain} gives
\begin{equation}\label{eq:complex-middle-A-interval}
 0\leq\omega^2\leq\frac13,\qquad
 \frac{2(1-\rho)}3\leq A\leq\frac{2-\rho}3.
\end{equation}
The left endpoint is precisely $a=0$, and the right endpoint is precisely
$a=b$.  Moreover, with
\[
 S_\omega(x)=\frac{\sin(\omega x)}{\omega},\qquad
 D_\omega(x)=e^x-\cos(\omega x)-S_\omega(x),
\]
one has the affine formula
\begin{equation}\label{eq:complex-middle-affine}
 K_{m,\omega}(x)=(1+\omega^2)S_\omega(x)+A D_\omega(x).
\end{equation}
For fixed $(x,\omega)$, the left side of
\eqref{eq:complex-root-middle-curvature} is a quadratic polynomial in $A$.
The following argument proves that it is nonnegative on the whole interval,
not merely at its endpoints.

\begin{proposition}[Complex-pair middle start]
\label{prop:complex-middle}
For every parameter pair in \eqref{eq:complex-middle-A-interval} and every
$x\geq0$, the stronger middle inequality
\eqref{eq:complex-root-middle-curvature} holds.
\end{proposition}

\begin{proof}
Put $z=\omega^2$ and $B=1-A+z$.  Then
\[
 K=Ae^x-A\cos(\omega x)+B\frac{\sin(\omega x)}{\omega}.
\]
Let $C_N$ be $N!$ times the coefficient of $x^N$ in
\[
 \mathcal C(x)=K^2+x^2(KK''-(K')^2).
\]
For $N\geq3$, set $M=N(N-1)/2$ and
\[
 R_N=(1+i\omega)^N
     +M(1-i\omega)^2(1+i\omega)^{N-2},
 \qquad J_N=\frac{\operatorname{Im}R_N}{\omega}.
\]
Writing the trigonometric part as a pair of conjugate exponentials gives the
exact identity
\begin{equation}\label{eq:middle-coefficient-pairing}
 C_N=A^22^N
 +2\operatorname{Re}\left[
   \left(\frac{-A-iB/\omega}{2}\right)^2(2i\omega)^N\right]
 -2A^2\operatorname{Re}R_N+2ABJ_N .
\end{equation}
All expressions at $\omega=0$ are interpreted by continuity.

The domain has the useful elementary consequences
\begin{equation}\label{eq:middle-domain-crude}
 0\leq z\leq A,\qquad 0\leq B\leq1,
 \qquad r:=\sqrt{1+z}\leq\frac2{\sqrt3}.
\end{equation}
Moreover,
\[
 J_N(0)=N+M(N-4)>0\quad(N\geq4).
\]
The polynomial $R_N(\omega)$ is a sum of a product of $N$ linear factors
and $M$ times another such product.  Taylor's theorem, using the conjugation
symmetry at the origin, therefore gives
\begin{equation}\label{eq:middle-J-remainder}
 J_N\geq J_N(0)-\frac z6(1+M)N(N-1)(N-2)r^{N-3}.
\end{equation}
Indeed, the absolute value of the third derivative of either product is at
most $N(N-1)(N-2)r^{N-3}$ on $[0,\omega]$.

We now use \eqref{eq:middle-J-remainder} in
\eqref{eq:middle-coefficient-pairing}, retain the nonnegative term
$2ABJ_N(0)$, and bound the other two real parts by their moduli.  Since
$z\leq A$ and $B\leq1$, for $N\geq6$ this yields
\begin{align}
 C_N\geq A^2\bigg\{2^N
 &-2(1+M)r^N
 -\frac13(1+M)N(N-1)(N-2)r^{N-3} \notag\\
 &-2^{N-1}\left(3^{-N/2}+3^{3-N/2}\right)\bigg\}.
 \label{eq:middle-tail-bound}
\end{align}
For example, the conjugate-pair term is bounded by
\[
 2^{N-1}\bigl(A^2z^{N/2}+B^2z^{N/2-1}\bigr)
 \leq
 2^{N-1}A^2\left(3^{-N/2}+3^{3-N/2}\right).
\]

The braces in \eqref{eq:middle-tail-bound} are positive for every
$N\geq25$.  After division by $2^N$, their total subtracted part is
\[
 E_N=\left(2(1+M)+14
 +\frac{\sqrt3}{8}(1+M)N(N-1)(N-2)\right)3^{-N/2}.
\]
At $N=25$ its numerator is
$616+519225\sqrt3$, whereas its denominator is $531441\sqrt3$, so
$E_{25}<1$.  Thereafter each of the two positive numerator summands grows
by a factor less than $4/3$, while the denominator grows by $\sqrt3$.
Indeed, their successive ratios are bounded by their values at $N=25$,
namely $666/616$ and $(326/301)(26/23)$, both less than $4/3$.
Thus $E_N$ decreases for $N\geq25$.

It remains only to check the finite range.  Parametrize the complete domain
by
\[
 \omega^2=\frac{1-\rho^2}{3},\qquad
 A=\frac{2-2\rho+\rho y}{3},
 \qquad 0\leq\rho,y\leq1.
\]
For $N=4,\ldots,24$, direct power-to-Bernstein conversion of $C_N$ gives
only nonnegative rational coefficients, first in $y$ and then in $\rho$.
This is a finite exact certificate: the derivative sequence used to form it
is
\[
 K^{(n)}(0)=A(1-c_nz^{\lfloor n/2\rfloor})
 +B s_nz^{\lfloor(n-1)/2\rfloor},
\]
where $(c_n)=(1,0,-1,0)$ and $(s_n)=(0,1,0,-1)$ periodically.  The short
exact-arithmetic verification is
\texttt{python/check\_complex\_middle\_boundaries.py}.  Finally,
$C_0=C_1=C_2=C_3=0$.  Hence every Taylor coefficient of $\mathcal C$ is
nonnegative, proving the result for $x\geq0$.
\end{proof}

Thus the middle start is now complete in both the real- and complex-spectrum
regions.  By the convolution argument preceding
\eqref{eq:middle-divided-difference}, the third start is complete as well.
At this stage, only the first canonical start remains open.

The calculation in \eqref{eq:real-first-simplex-curvature} in fact has a
root-free extension which proves a substantial part of the remaining case.
Assume $Q>0$ and define
\[
 \mathcal M(t)=\frac{h_1(t)-1}{Qt^3}.
\]
Since $h_1'=c h_3$, the function $\mathcal M$ is, up to a positive constant,
$t^{-3}\int_0^t h_3(s)\,ds$.  The completed third-start result and the
convolution lemma imply that $\mathcal M$ is log-convex.  Its series at the
origin gives
\[
 (\log\mathcal M)'(0)=\frac{a+b}{4}\geq0;
\]
hence $(\log\mathcal M)'\geq0$ everywhere.  Put
\[
 \vartheta=\frac{h_1-1}{h_1},\qquad
 \mathcal L=3+t(\log\mathcal M)',\qquad
 \mathcal V=t^2(\log\mathcal M)''.
\]
Then $\mathcal L\geq3$, $\mathcal V\geq0$, and direct differentiation gives
\begin{equation}\label{eq:first-integrated-third-identity}
 -t^2(\log h_1)''
 =\vartheta\bigl(3-\mathcal V
 -(1-\vartheta)\mathcal L^2\bigr).
\end{equation}
It follows immediately that the first-start target holds, for every spectrum,
throughout the region
\begin{equation}\label{eq:first-small-return-region}
 h_1(t)-1\leq5+3\sqrt3.
\end{equation}
Indeed, this is equivalent to
$\vartheta\leq(1+\sqrt3)/3$, and the right side of
\eqref{eq:first-integrated-third-identity} is then at most
$\vartheta(9\vartheta-6)\leq2$.  Outside
\eqref{eq:first-small-return-region}, the sole remaining scalar inequality is
\begin{equation}\label{eq:first-integrated-third-target}
 \mathcal V+(1-\vartheta)\mathcal L^2
 \geq3-\frac2\vartheta.
\end{equation}
For real roots, $\mathcal V$ is precisely the tilted-simplex variance $V$ in
\eqref{eq:real-first-simplex-target}; for a complex pair it remains a
nonnegative convolution curvature.  Thus
\eqref{eq:first-integrated-third-target} unifies the two spectral regions.

Exact computation supports the cleaner, stronger inequality
\begin{equation}\label{eq:first-integrated-third-strong}
 \mathcal V+(1-\vartheta)\mathcal L^2\geq\vartheta.
\end{equation}
It implies \eqref{eq:first-integrated-third-target}, since
$\vartheta\geq3-2/\vartheta$ for $0<\vartheta\leq1$.  In the original
variables it is equivalent to
\begin{equation}\label{eq:first-strong-curvature}
 t^2\bigl(h_1h_1''-(h_1')^2\bigr)+2h_1^2-h_1-1\geq0,
\end{equation}
and would give the sharper bound
$E_1\leq\vartheta(3-\vartheta)\leq2$.
At coefficient level, \eqref{eq:first-strong-curvature} replaces
$\Gamma(C,J,K)$ by $\Gamma(C,J,K)-\alpha(C,J,K)$.  Exact tests find the same
five negative triples $(2,0,K)$, $3\leq K\leq7$, and no others.  Grouping
with the neighboring $C=1,3$ terms gives positive quadratics whose
discriminants are
\[
 -5907200,\quad -1402143,\quad -7154704640,\quad
 -11951446784,\quad -263131073024.
\]
Thus the two normalized coefficient monotonicities below would prove even
the strengthened form \eqref{eq:first-strong-curvature}; no additional
infinite family is introduced.

The completed third bound almost closes the first one by a barrier argument.
This gives a useful quantitative target for strengthening the convolution
step.  Introduce logarithmic time $\tau=\log t$ and the scaled variables
\[
 Z=tz,\qquad P=tp,\qquad R=tq,
 \qquad A_0=ta,\qquad B_0=tb.
\]
Then \eqref{eq:canonical-rate-system} becomes
\[
 \dot Z=Z(1+B_0+R-Z),\qquad
 \dot P=P(1+Z-A_0-P),\qquad
 \dot R=R(1+A_0+P-B_0-R),
\]
where a dot denotes $d/d\tau$.  Set
\[
 E_1=Z(Z-B_0-R),\qquad
 E_3=R(B_0+R-A_0-P).
\]
A direct differentiation, followed only by the substitution $E_1=2$, gives
the exact first-contact identity
\begin{equation}\label{eq:first-contact-slack}
 \left.\dot E_1\right|_{E_1=2}
 =Z(E_3-2)+4-\frac4Z.
\end{equation}
At such a contact one has $Z>\sqrt2$.  The established third inequality
$E_3\leq2$ alone therefore falls just short of forcing $\dot E_1\leq0$.
It would be enough to prove the following quantitative slack implication:
\begin{equation}\label{eq:first-contact-slack-target}
 E_1=2\quad\Longrightarrow\quad
 2-E_3\geq\frac{4(Z-1)}{Z^2}.
\end{equation}
The left side is exactly
$t^2(\log(t^{-2}h_3))''$, so
\eqref{eq:first-contact-slack-target} asks for a quantitative version of the
same convolution log-convexity argument that already proves its
nonnegativity.

That slack itself has an exact positive decomposition.  Put
\[
 F(s)=\frac{e^{-bs}h_2(s)}s,
 \qquad
 G(t)=\frac{e^{-bt}h_3(t)}{t^2}
     =x_3\int_0^1uF(tu)\,du,
\]
and let $\mu_t$ be the probability measure on $[0,1]$ with density
proportional to $uF(tu)$.  If
\[
 E_2(s)=s^2p(s)(a+p(s)-z(s)),
 \qquad L(s)=s(\log F)'(s),
\]
then differentiation under the integral gives
\begin{equation}\label{eq:third-slack-decomposition}
 2-E_3(t)
 =\E_{\mu_t}[1-E_2(tu)]
  +\operatorname{Var}_{\mu_t}(L(tu)).
\end{equation}
Both terms are nonnegative by the stronger middle bound.  Combining
\eqref{eq:first-contact-slack-target} and
\eqref{eq:third-slack-decomposition} isolates the remaining dynamic question:
at a first contact, one must extract a lower bound on either the averaged
middle slack or the spread of the logarithmic slopes.  No oscillatory root
variables occur in this formulation.

There is a complementary coefficient reduction for the first canonical
start.  It sharply localizes the cancellation that a proof still has to
explain.  Write
\[
 S=a+b,\qquad \Delta=(a-b)^2,\qquad Q=cx_2x_3.
\]
The ordinary generating function of the return-walk weights is
\begin{align}
 \sum_{n\geq0}(B^n)_{11}z^n
 &=\frac{(1-az)(1-bz)}
         {(1-az)(1-bz)-Qz^3} \notag\\
 &=\sum_{C\geq0}
   \frac{Q^Cz^{3C}}
        {\bigl((1-Sz/2)^2-\Delta z^2/4\bigr)^C}.
 \label{eq:canonical-positive-coordinates}
\end{align}
Here the term with $C=0$ is $1$.  For $C\geq1$, expansion of the last
display gives the manifestly nonnegative coefficient
\begin{equation}\label{eq:canonical-alpha}
 \alpha(C,J,K)=
 \binom{C+J-1}{J}
 \binom{2C+2J+K-1}{K}\,2^{-2J-K}
\end{equation}
of
$Q^C\Delta^JS^Kz^{3C+2J+K}$.

This turns coefficientwise curvature into a finite binomial-sum question.
Set $\alpha(0,0,0)=1$, with all other $\alpha(0,J,K)=0$, put
$N=3C+2J+K$, and define
\begin{equation}\label{eq:canonical-gamma}
 \begin{aligned}
 \Gamma(C,J,K)
 &=\sum_{\substack{0\leq c\leq C\\0\leq j\leq J\\0\leq k\leq K}}
 \binom{N}{3c+2j+k}
 \left(2-\frac N2+
       \frac{(N-2(3c+2j+k))^2}{2}\right)\\
 &\hspace{34mm}{}
 \times\alpha(c,j,k)\alpha(C-c,J-j,K-k).
 \end{aligned}
\end{equation}
Thus $\Gamma(C,J,K)$ is exactly the coefficient of
$Q^C\Delta^JS^K$ in the $N$th exponential Taylor coefficient of the
corrected curvature.

There is a useful probabilistic simplification of
\eqref{eq:canonical-gamma}.  Let
\[
 T(C,J,K)=\sum_{c,j,k}W(c,j,k),
 \qquad
 \pi(c,j,k)=\frac{W(c,j,k)}{T(C,J,K)},
\]
with $W$ as in \eqref{eq:S-loop-insertion} below, and regard
$n=3c+2j+k$ as a random variable under $\pi$.  Complementing the cut sends
$n$ to $N-n$ without changing its weight, so $\E_\pi n=N/2$.  Since the
kernel in \eqref{eq:canonical-gamma} is
\[
 2-\frac N2+2\left(n-\frac N2\right)^2,
\]
one obtains the exact identity
\begin{equation}\label{eq:canonical-cut-variance}
 \Gamma(C,J,K)=T(C,J,K)
 \left(2-\frac N2+2\operatorname{Var}_\pi(n)\right).
\end{equation}
Thus the uniform coefficient sign assertion is equivalent to the elementary
cut-dispersion bound
\begin{equation}\label{eq:canonical-cut-variance-target}
 \operatorname{Var}_\pi(n)\geq\frac N4-1.
\end{equation}
The five exceptional triples below are exactly the cases in which this
conditional cut distribution is slightly too concentrated.  This variance
formulation is also compatible with conditioning first on the number $c$ of
cycles to the left of the cut, so the law of total variance offers a second
route to the desired finite-exception statement.

Exact expansion suggests the following particularly rigid sign pattern:
\[
 \Gamma(C,J,K)\geq0
\]
except when $(C,J)=(2,0)$ and $3\leq K\leq7$.  This has been checked in
exact integer arithmetic for
$1\leq C\leq12$, $0\leq J\leq9$, and $0\leq K\leq20$; it is not yet a
proof for all indices.  The five exceptions do not threaten the desired
curvature.  They occur at total degrees $N=9,\ldots,13$ and combine with
the neighboring $C=1$ and $C=3$ terms into positive quadratics.  After
discarding other nonnegative terms, the relevant brackets are
\[
\begin{array}{c|c}
N&\text{quadratic in }x=Q/S^3\\ \hline
9&133-320x+6592x^2\\
10&47-437x+18096x^2\\
11&171-3280x+309952x^2\\
12&85-2296x+569728x^2\\
13&55-920x+1179744x^2.
\end{array}
\]
Their discriminants are respectively
\[
 -3404544,\quad -3211079,\quad -201248768,\quad
 -188435904,\quad -258697280,
\]
so all five are strictly positive.  Consequently a proof of the uniform
sign assertion for \eqref{eq:canonical-gamma} would settle the first
canonical start coefficient by coefficient.  Analogous expansions for the
other two starts show no exceptional negative group in the tested range.
This finite-exception binomial lemma is now the most concrete algebraic
target alongside the Riccati inequalities.

There is evidence for a sharper reduction of that lemma.  Normalize by the
positive coefficient in \eqref{eq:canonical-alpha} and set
\[
 R(C,J,K)=\frac{\Gamma(C,J,K)}{\alpha(C,J,K)}.
\]
Exact arithmetic through $C\leq9$, $J\leq14$, and $K\leq20$ gives
\begin{align*}
 R(C,J+1,K)&\geq R(C,J,K),\qquad &&C\geq1,\\
 R(C,J,K+1)&\geq R(C,J,K),
 &&(C,J,K)\notin\{(2,0,k):0\leq k\leq5\}.
\end{align*}
If these two monotonicity statements hold in general, the whole uniform sign
claim reduces to the already proved equal-rate value
\[
 R(C,0,0)=
 \frac{2^{3C+1}+(-1)^C(-27C^2+9C+4)}{3}
\]
and the single slice $C=2$, $J=0$.  That slice has the closed form
\begin{equation}\label{eq:C2-closed-form}
 \begin{aligned}
 \Gamma(2,0,K)=\frac{1}{6\,2^K}\bigl\{
 &96\,2^K(K^2-5K-34)\\
 &{}+K^5+23K^4+219K^3+1081K^2+2888K+3348
 \bigr\}.
 \end{aligned}
\end{equation}
Formula~\eqref{eq:C2-closed-form} follows by writing the interior summation
with $n=k+3$, extending it to the full binomial sum $0\leq n\leq K+6$,
evaluating the moments through degree four, and subtracting the six boundary
terms.
It is positive for $K\geq9$ directly, because both displayed groups are then
nonnegative, and inspection of $0\leq K\leq8$ leaves exactly
$K=3,\ldots,7$ negative.  Thus the two normalized monotonicity inequalities,
rather than an unrestricted three-variable sign calculation, are sufficient
for the first start.

There is an exact insertion formula for the second monotonicity statement.
It substantially narrows what remains to be proved.  Put
\[
 \kappa_N(n)=2-\frac N2+\frac{(N-2n)^2}{2},\qquad
 W(c,j,k)=\binom Nn\alpha(c,j,k)
 \alpha(C-c,J-j,K-k),
\]
where $n=3c+2j+k$ and $N=3C+2J+K$.  Inserting one $S$-loop on either side
of the cut gives
\begin{equation}\label{eq:S-loop-insertion}
 R(C,J,K+1)-R(C,J,K)
 =\frac1{\alpha(C,J,K)}\sum_{c,j,k}W(c,j,k)\Phi_{N,C}(n,c),
\end{equation}
where
\[
 \Phi_{N,C}(n,c)
 =\frac{N+1}{N-C}\left\{
 \frac{n-c}{n+1}\kappa_{N+1}(n+1)
 +\frac{N-n-C+c}{N+1-n}\kappa_{N+1}(n)
 \right\}-\kappa_N(n).
\]
The identity uses only
\[
 (K+1)\alpha(C,J,K+1)=\frac{N-C}{2}\alpha(C,J,K)
\]
and its two cut analogues.  In the centered variables
$u=N-2n$ and $r=C-2c$, the kernel becomes
\begin{equation}\label{eq:S-loop-centered-kernel}
 \Phi_{N,C}(n,c)=
 \frac{A_N(u^2)}{2((N+2)^2-u^2)}
 +\frac{(N+1)(3N-u^2)u(u-r)}
 {(N-C)((N+2)^2-u^2)},
\end{equation}
with
\[
 A_N(v)=v^2+(N^2-3N)v-N^3+2N^2+8N.
\]
Both $W$ and $\Phi$ are invariant under the complementary-cut involution
$(n,c)\mapsto(N-n,C-c)$.  The kernel can be negative, so pointwise insertion
does not prove monotonicity.  Equation~\eqref{eq:S-loop-insertion} shows that
the missing statement is precisely a correlation estimate between the total
cut imbalance $u$ and its noncycle part $u-r$; it replaces a comparison of
two triple sums by one weighted average over the old cuts.

The first normalized monotonicity has a parallel exact insertion formula.
It is useful because it puts both conjectural comparisons on the same old-cut
probability space.  With $N,n,W$ as above, set
\[
 m=n-c,\qquad \bar m=N-n-C+c,
\]
and define
\begin{align*}
 \lambda_L&=\frac{(N+2)(N+1)m(m+1)}
 {8(n+2)(n+1)(2c+2j+1)},\\
 \lambda_R&=\frac{(N+2)(N+1)\bar m(\bar m+1)}
 {8(N-n+2)(N-n+1)(2(C-c)+2(J-j)+1)},\\
 \Psi_{C,J,K}(c,j,k)
 &=\frac{8(2C+2J+1)}{(N-C)(N-C+1)}
 \left\{\lambda_L\kappa_{N+2}(n+2)
       +\lambda_R\kappa_{N+2}(n)\right\}-\kappa_N(n).
\end{align*}
Then
\begin{equation}\label{eq:Delta-pair-insertion}
 R(C,J+1,K)-R(C,J,K)
 =\frac1{\alpha(C,J,K)}
   \sum_{c,j,k}W(c,j,k)\Psi_{C,J,K}(c,j,k).
\end{equation}
Indeed,
\[
 \frac{\alpha(c,j+1,k)}{\alpha(c,j,k)}
 =\frac{(n-c)(n-c+1)}
 {8(j+1)(2c+2j+1)},
\]
and the same identity on the complementary side accounts for every newly
inserted pair exactly $J+1$ times.  Formula
\eqref{eq:Delta-pair-insertion} follows after using the corresponding total
ratio and the two binomial-coefficient ratios.  Like the loop kernel,
$\Psi$ can be negative pointwise; the assertion needed is its nonnegative
weighted mean.

There is also a positive beta-integral form which eliminates the loop split
$k$ from either weighted average.  Put
\[
 \beta(C,J,K)=\frac{\alpha(C,J,K)}{(3C+2J+K)!}.
\]
For $C\geq1$, a beta integral gives
\begin{equation}\label{eq:canonical-beta-alpha}
 \beta(C,J,K)=\frac{d_{C,J}}{2^K K!}
 \int_0^1x^{2C+2J+K-1}(1-x)^C\,dx,
 \qquad
 d_{C,J}=\frac{(C)_J}
 {4^J J!(2C+2J-1)!C!}.
\end{equation}
For an interior split $1\leq c\leq C-1$, fix $j$ and introduce independent
integration variables $x,y$ from the two copies of
\eqref{eq:canonical-beta-alpha}.  If
\[
 p=\frac{x}{x+y},\qquad d=3c+2j,
\]
then summing all $k=0,\ldots,K$ replaces the loop allocation by
$X\sim\operatorname{Bin}(K,p)$ and gives
\begin{align}
 &\sum_{k=0}^K
 \kappa_N(d+k)\beta(c,j,k)\beta(C-c,J-j,K-k) \notag\\
 &\quad=\frac{d_{c,j}d_{C-c,J-j}}{2^K K!}
 \int_0^1\!\int_0^1
 x^{2c+2j-1}(1-x)^c
 y^{2(C-c)+2(J-j)-1}(1-y)^{C-c} \notag\\
 &\hspace{42mm}{}
 \times(x+y)^K\,\Xi_{N,K,d}(p)\,dx\,dy,
 \label{eq:canonical-beta-cut}
\end{align}
where the remaining kernel is merely the quadratic
\begin{equation}\label{eq:canonical-beta-kernel}
 \Xi_{N,K,d}(p)=2-\frac N2
 +2\left(d+Kp-\frac N2\right)^2+2Kp(1-p).
\end{equation}
The two boundary cuts contribute
$2\kappa_N(0)\beta(C,J,K)$.  Since
$\Gamma/N!=\sum\kappa_N(n)\beta(c,j,k)
\beta(C-c,J-j,K-k)$, equations
\eqref{eq:canonical-beta-alpha}--\eqref{eq:canonical-beta-kernel} are an
exact double-integral formula for $R(C,J,K)$.  Thus the $k$-sum and its
changing support are not intrinsic: after conditioning on $x,y$, loop
allocation is an ordinary binomial random variable.  What remains is to
control the negative constant in \eqref{eq:canonical-beta-kernel} by the
square, binomial variance, and the two positive boundary cuts.  This form
uses the actual $\alpha$-weights and therefore avoids the false generic
renewal-point variance statement, which ignores the averaging over loop
allocations.

For exact computation there is an even simpler coefficient system which
keeps the rate ordering explicit.  Put $d=b-a\geq0$.  Then the coefficient
of
\[
 Q^Ca^Jd^Kz^{3C+J+K}
\]
in the first-start return series is
\begin{equation}\label{eq:canonical-ordered-alpha}
 \widetilde\alpha(C,J,K)=
 \binom{C+K-1}{K}\binom{2C+K+J-1}{J}
 \qquad(C\geq1).
\end{equation}
To see this, expand
$(1-az)^{-C}(1-(a+d)z)^{-C}$ and apply Vandermonde's identity.
The analogous coefficients for the second and third starts are respectively
\[
 \binom{C+K-1}{K}\binom{2C+K+J}{J},
 \qquad
 \binom{C+K}{K}\binom{2C+K+J+1}{J},
\]
with the evident convention at $C=0$.  In exact tests through total Taylor
degree $20$, the corrected-curvature coefficients for starts two and three
are all nonnegative.  For start one the only negative monomials again have
$C=2$ and total Taylor degree $9$ through $13$: in the ordered coordinates
they are the two families $Q^2a^Ld^0$ and $Q^2a^{L-1}d$,
$3\leq L\leq7$.  Formula~\eqref{eq:canonical-ordered-alpha} is often more
convenient than \eqref{eq:canonical-alpha} for a direct Vandermonde proof;
the symmetric coordinates remain better for combining the exceptional
terms into the five quadratics above.

The boundary $a=0$ is especially simple in these coordinates.  Write
$\widetilde\Gamma(C,K)$ for the coefficient of $Q^Cd^K$ in the corrected
first-start curvature and put
\[
 \widetilde R(C,K)=
 \frac{\widetilde\Gamma(C,K)}{\widetilde\alpha(C,0,K)},
 \qquad N=3C+K.
\]
There is an exact loop-insertion identity parallel to
\eqref{eq:S-loop-insertion}.  If
\[
 \widetilde W(c,k)=\binom N{3c+k}
 \widetilde\alpha(c,0,k)
 \widetilde\alpha(C-c,0,K-k),
 \qquad n=3c+k,
\]
then
\begin{equation}\label{eq:ordered-loop-insertion}
 \widetilde R(C,K+1)-\widetilde R(C,K)
 =\frac1{\widetilde\alpha(C,0,K)}
  \sum_{c,k}\widetilde W(c,k)\widetilde\Phi_{C,K}(c,k),
\end{equation}
where
\begin{align*}
 \widetilde\Phi_{C,K}(c,k)
 ={}&\frac{N+1}{C+K}\left\{
 \frac{c+k}{n+1}\kappa_{N+1}(n+1)
 +\frac{C-c+K-k}{N+1-n}\kappa_{N+1}(n)
 \right\}-\kappa_N(n).
\end{align*}
Indeed, mark and remove one of the $K+1$ loops in the new coefficient, on
either side of the cut, and use
\[
 (K+1)\widetilde\alpha(C,0,K+1)
 =(C+K)\widetilde\alpha(C,0,K).
\]
This accounts for every new term exactly $K+1$ times and proves
\eqref{eq:ordered-loop-insertion}.

The initial value $\widetilde R(C,0)$ is the proved equal-rate expression
in \eqref{eq:equal-rate-coefficient}.  Exact arithmetic finds
\[
 \widetilde R(C,K+1)\geq\widetilde R(C,K)
\]
for $1\leq C\leq12$ and $0\leq K\leq20$, with the sole exception
$(C,K)=(2,0)$.  That exception is harmless, and in fact the entire
$C=2$ slice can be evaluated explicitly:
\begin{equation}\label{eq:ordered-C2-closed-form}
 \widetilde\Gamma(2,K)
 =128\,2^K-\frac{K^4+14K^3+69K^2+156K+228}{2}>0.
\end{equation}
To verify the strict inequality, check $K=0,1,2$ directly.  Thereafter it
propagates by induction, since for
$P(K)=K^4+14K^3+69K^2+156K+228$ one has
\[
 2P(K)-P(K+1)=K^4+10K^3+21K^2-28K-12>0
 \qquad(K\geq2).
\]
The $C=1$ slice is even simpler:
$\widetilde\Gamma(1,K)=K^2+5K+10$.
Consequently, proving that the weighted mean in
\eqref{eq:ordered-loop-insertion} is nonnegative for $C\geq3$ would give an
independent coefficientwise proof of the $a=0$ boundary already established
in Proposition~\ref{prop:first-a-zero}.  The kernel itself changes sign, so
this assertion is again a correlation inequality rather than a termwise
estimate.  All identities and the quoted finite check are verified
in \path{python/check_canonical_k3_coefficients.py}.

\section{Equivalent walk-count formulation}

The Taylor expansion
\[
  F(t)=\sum_{n\geq0}\frac{t^n}{n!}\,u^TA^nv
\]
has nonnegative coefficients.  Define a probability distribution by
\[
  \Pr_t(N=n)=
  \frac{t^n u^TA^nv}{n!F(t)}.
\]
Then
\[
  \E_tN=t\frac{F'(t)}{F(t)},
  \qquad
  \Var_t(N)=t\frac{F'(t)}{F(t)}
             +t^2(\log F(t))''.
\]
Consequently \eqref{eq:matrix-curvature} is exactly
\begin{equation}\label{eq:walkvariance}
  \Var_t(N)\geq \E_tN-(s-1).
\end{equation}
Here \(u^TA^nv\) is the total weight of length-\(n\) walks in a weighted directed
graph with \(s\) vertices.  Thus the matrix lemma says that conditioning a
Poisson walk by its endpoints can reduce its variance below its mean by at most
\(s-1\).  A likely combinatorial proof should remove loops from a walk: at most
\(s-1\) steps survive loop erasure, while all remaining steps must carry the
Poisson-type variance on the right side of \eqref{eq:walkvariance}.

It is tempting to seek a coefficientwise strengthening of this formulation.
Write
\[
 D(t)=t^2\bigl(F(t)F''(t)-F'(t)^2\bigr)+(s-1)F(t)^2
     =\sum_{N\geq0}D_N\frac{t^N}{N!},
\]
and put \(a_n=u^TA^nv\).  Direct expansion and symmetrization give
\begin{equation}\label{eq:coefficient-formula}
 D_N=\sum_{i=0}^N\binom Ni
 \left(s-1-\frac N2+\frac{(N-2i)^2}{2}\right)a_i a_{N-i}.
\end{equation}
If \(I\) has the symmetric distribution
\[
 \Pr(I=i)=
 \frac{\binom Ni a_i a_{N-i}}
      {\sum_{j=0}^N\binom Nj a_j a_{N-j}},
\]
then \(D_N\geq0\) is precisely
\begin{equation}\label{eq:split-variance}
 \Var(I)\geq\frac N4-\frac{s-1}{2}.
\end{equation}
The normalizing denominator counts two weighted walks of total length \(N\),
together with an interleaving of their steps.  Thus \(I\) is the length of
the first walk.  However, the proposed coefficientwise inequality is false,
even though all numerical tests of the pointwise sum \(D(t)\) remain
nonnegative.  Take \(s=4\), let \(A\) be the adjacency matrix of the directed
four-cycle (all four edge weights equal to one), and take \(u=v=e_1\).  Then
\[
 a_n=\begin{cases}1,&4\mid n,\\0,&4\nmid n.\end{cases}
\]
For \(N=8\), only \(i=0,4,8\) contribute to
\eqref{eq:coefficient-formula}, and hence
\[
 D_8=31+31-\binom84=-8<0.
\]
Equivalently, \(I\) is supported on \(0,4,8\), with weights \(1,70,1\), so
\[
 \Var(I)=\frac49<\frac12
       =\frac N4-\frac{s-1}{2}.
\]
Thus cancellation between different Taylor coefficients is essential; a
coefficient-by-coefficient tail-switching proof cannot establish
Conjecture~\ref{conj:matrix}.  Notice that the same example gives
\[
 F(t)=\frac{\cosh t+\cos t}{2},
\]
and it does not furnish a pointwise counterexample to
\eqref{eq:matrix-curvature}.
\section{A state-return route beyond three moments}

There is a second possible route to the all-state conjecture which does not
pass through the matrix-exponential statement.  Remove one transient state
and decompose the absorption time into a geometric number of return
excursions followed by one terminal excursion.  If \(x\) is the mean holding
time at the removed state and \(r\) is the mean number of returns, the moment
generating functions satisfy
\[
 M_Z(z)=\frac{V(z)}{(1+r)(1-xz)-rU(z)}.
\]
This suggests the following uniform induction lemma: if both excursion laws,
of order at most \(n-1\), satisfy the hierarchy through level \(k\), then the
renewal law satisfies level \(k\) at order \(n\).  Such a lemma for every \(k\)
would prove the full conjecture by induction on the number of states.  It is
already known at \(k=1,2,3\).

The first unproved instance, \(k=4\), now has a reasonably small exact
reduction.  Write the first five return-excursion moments as
\[
 (u,A,C,E,G)
\]
and the terminal-excursion moments as \((v,B,D,F,J)\).  After adjoining the
exponential holding time, let \(a_j,b_j\) be the return and terminal kernel
moments.  Then
\[
 M_j=b_j+ra_j+\sum_{\ell=1}^{j-1}\binom j\ell ra_\ell M_{j-\ell}.
\]
The sharp fifth-moment child bounds permit the substitutions
\[
 G=\frac{n+3}{n+2}\frac{E^2}{C},
 \qquad
 J=\frac{n+3}{n+2}\frac{F^2}{D}.
\]
The resulting parent gap
\[
 \mathcal G_4=(n+3)M_3M_5-(n+4)M_4^2
\]
is a strictly convex quadratic in \(E,F\).  If

\[
 R_3=M_3-rC-D,
\]
its Hessian has off-diagonal entry \(-2r(n+4)\), and its determinant is
\[
 \frac{4r(n+3)^2M_3}{CD(n+2)^2}
 \left(rC+D+(n+3)^2R_3\right)>0.
\]
Thus the top-moment convexity used in the \(k=3\) proof persists at \(k=4\).

The stationary-point algebra simplifies substantially after a weighted
change of variables.  Put
\[
 \alpha=\frac{n+3}{n+2},\quad
 c=rC,\quad d=D,\quad w=c+d,\quad
 y=\frac EC,\quad f=\frac FD,
\]
and set
\[
 m=M_1,\qquad P=M_3,\qquad T=x+ra_1,
\]
\[
 L=M_4\big|_{E=F=0},\qquad
 K=M_5\big|_{E=F=0}.
\]
Then the entire quadratic is
\begin{align}
 \mathcal G_4={}&(n+3)P
 \{K+\alpha(cy^2+df^2)+5c(T+m)y+5dTf\}\notag\\
 &\quad -(n+4)(L+cy+df)^2.\label{eq:k4-top-quadratic}
\end{align}
Introduce the weighted mean and difference
\[
 \mu=\frac{cy+df}{w},\qquad \delta=f-y.
\]
The two variables separate in \eqref{eq:k4-top-quadratic}.  In particular,
every interior stationary point satisfies the simple identity
\begin{equation}\label{eq:k4-stationary-difference}
 \delta_*=\frac{5m}{2\alpha}
 =\frac{5(n+2)}{2(n+3)}M_1.
\end{equation}
After minimizing in \(\delta\), the remaining quadratic in \(\mu\) has leading
coefficient
\[
 A_\mu=\frac{w}{n+2}\Delta,
 \qquad
 \Delta=(n+3)^2(P-w)+w>0,
\]
linear coefficient
\[
 B_\mu=5(n+3)P(wT+cm)-2(n+4)Lw,
\]
and constant term
\[
 C_\mu=(n+3)P\left(K-\frac{25cdm^2}{4\alpha w}\right)
       -(n+4)L^2.
\]
Consequently the value at the unconstrained interior stationary point is
\begin{equation}\label{eq:k4-interior-value}
 C_\mu-\frac{B_\mu^2}{4A_\mu}.
\end{equation}
The lower fourth-moment child bounds say
\[
 y\geq\frac{n+2}{n+1}\frac CA,
 \qquad
 f\geq\frac{n+2}{n+1}\frac DB.
\]
The algebraic lemma established below is that
\eqref{eq:k4-interior-value} is nonnegative whenever the stationary point
given by \eqref{eq:k4-stationary-difference} satisfies these two feasibility
inequalities and the lower child hierarchy holds.

The quantities in this last lemma also obey a short recurrence.  With
\(a=rA\) and \(b=B\),
\begin{align*}
 M_2&=2Tm+a+b,\\
 P&=3TM_2+3am+c+d,\\
 L&=4TP+6aM_2+4cm,\\
 K&=5TL+10aP+10cM_2.
\end{align*}
These identities remove all large kernel expansions from the interior
problem.  Random searches over the child moment cone have found feasible
interior stationary points, but none with a negative value in
\eqref{eq:k4-interior-value}.

There is a sharper reduction of this last problem.  By homogeneity put
\(m=1\), write \(f=y+5(n+2)/(2(n+3))\), and let
\(p,q,A_0,B_0\in[0,1]\).  The feasibility inequalities and the two
third-moment child inequalities are parametrized by
\[
 a=\frac{nTypA_0}{n+2},\qquad
 c=\frac{ayp}{\beta},\qquad
 b=\frac{n(1-T)fqB_0}{n+2},\qquad
 d=\frac{bfq}{\beta},
 \qquad \beta=\frac{n+2}{n+1}.
\]
Indeed \(p,q\leq1\) are exactly the two stationary-point feasibility
bounds, while \(A_0,B_0\leq1\) encode
\(nTc\geq(n+1)a^2\) and
\(n(1-T)d\geq(n+1)b^2\).  The remaining terminal second-moment bound is
\[
             \frac{(n-1)B_0qf}{n+2}\geq1-T.
\]
Define
\[
\begin{split}
 S&=2(n+4)M_4-(n+3)P(2\alpha f+5T),\\
 R&=M_5-M_4\left(\alpha f+\frac52T\right).
\end{split}
\]
Then the parent gap has the exact identity
\[
 \mathcal G_4=(n+3)PR-\frac12M_4S.
\]
Both \(S\) and \(R\) are affine in \(B_0\).  If
\[
 S=S_0+B_0S_1,\qquad R=R_0+B_0R_1,\qquad
 J=R_0S_1-R_1S_0,
\]
then at an interior stationary point \(S=0\),
\[
                         R=\frac{J}{S_1}.
\]
On finite positive parameters, the exact certificates below prove
\[
             S(1)\geq0\quad\Longrightarrow\quad S_1>0,\ J\geq0.
\]
The strict slope and the orientation argument below justify both division
by $S_1$ and the condition $S(1)\geq0$ at a stationary point. The terminal
second-moment bound is unnecessary for this interior step.

After compactifying \(n-4\) and \(y\), the sole common infinity corner of
\(S(1),S_1,J\) has Newton boundary
\[
                         (0,1),\qquad(2,0).
\]
Thus it is covered by the two charts associated with
\(1/y\asymp1/n^2\).  Their radial faces reduce the remaining implication to
five-variable polynomials with respectively \(14,4,20\) terms.  On one
secondary face the slope implication follows directly from convexity in
\(A_0\): if
\[
 H=12A_0p-2\theta q+3q-6,
\]
then the condition polynomial is convex in \(A_0\), is nonpositive at
\(A_0=0\), and at the root of \(H\) contains the nonnegative factor
\[
                6-3p+T(3-2\theta)(q-p).
\]
Hence the condition being nonnegative forces \(H\geq0\).

The analogous slope implication on the opposite infinity face is now also
exact.  Write \(F=S(1)/4\), \(H=S_1/2\), and put
\[
 K_0=3TWq-6TW-2q,\qquad
 L=3TWp-3TWq+3Wp-6W-2p+2q.
\]
Then \(H=12A_0TWp+K_0\), \(F(0)=-q(T-1)K_0\leq0\), and at the root of
\(H\),
\[
                         F=-\frac{K_0L}{12W}.
\]
Convexity in \(A_0\) therefore leaves only endpoint estimates.  With
\(C=3W(T+1)-2\) and \(D=2-3TW\), the constraint \(H(1)\geq0\) gives
\(p\geq\tfrac12+Dq/(12TW)\).  If \(C\geq0\), maximizing
\(L=Cp+Dq-6W\) at the appropriate endpoints gives either \(-3W\) or
\(3W(T-1)-2\).  If \(C\leq0\), substitution of this lower bound leaves
\(-Q/(12TW)\); \(Q\) is affine in \(q\), its \(q=0\) endpoint is
manifestly nonnegative, and its \(q=1\) endpoint has nonnegative Bernstein
coefficients after \(W=2z/[3(T+1)]\) and one split at \(z=1/2\).
For the remaining endpoint put \(E=15TW+3W-2\).  When \(E\leq0\), the
polynomial \(-3(5T+1)F(1)\), after
\(W=2z/[3(5T+1)]\), has nonnegative Bernstein coefficients.  When
\(E\geq0\), convexity in \(p\) reduces to the preceding root or to
\(p=1\); in the latter case \(H\leq0\) gives
\(q\geq6TW/(2-3TW)\), and \(F(1)\) is decreasing in \(q\) and already
nonpositive at that lower endpoint.  Thus \(S(1)\geq0\) implies
\(S_1\geq0\) on both infinity faces.

One resultant face is exact as well.  On the \(T\leq Z\) secondary chart
(which is the same polynomial chart as the opposite primary face), put
\(K=J/(8A_0p)\).  An exact rational identity has been found of the form
\[
                         K=M F+N H+P,
\]
where every Bernstein coefficient of \(M,N,P\) is nonnegative.  Hence
\(F\geq0\) implies \(J\geq0\) there.

On the remaining \(Z\leq T\) chart, both \(F\) and
\(K=J/(8A_0p)\) are affine in \(T\).  The endpoint \(T=0\) now has another
exact identity \(K=MF+NH+P\) with nonnegative rational Bernstein
coefficients.  At \(T=1\), put \(x=A_0p\).  Then
\[
 F=2x\{6x+(3-R)p-6\}.
\]
At its boundary \(x_0=1-(3-R)p/6\), the target factors as
\[
 K(x_0)=\frac13
 \{2Rp-2Rq-6p+3q+6\}
 \{10Rp^2-3Rp-30p^2+60p-18\}.
\]
Feasibility gives \(p\geq6/(9-R)\), which makes both factors
nonnegative.  If \(p\geq q\), the target is convex in \(x\) and its
derivative at \(x_0\) is nonnegative; if \(p\leq q\), it is concave and
its second endpoint \(x=p\) is nonnegative by checking the two affine
endpoints \(q=p,1\).  Thus both \(T\)-endpoints are exact.  The interior crossing
\(F(T)=0\) is now exact as well.  Its unresolved boxes converge to the
single equality corner
\[
 T=1,\qquad R=0,\qquad p=1,\qquad q=0,\qquad
 x=A_0p=\frac12,\qquad H=0.
\]
Put \(u=1-T\), \(a=1-p\), and \(h=H\), and use the weighted variables
\[
                         h,\qquad 6a,\qquad3q,\qquad2R,\qquad3u.
\]
In the chart where \(h\) dominates, the exact Bernstein coefficients of
\(K/h\) all lie between $2$ and $10$.  In each of the other four
charts, all Bernstein coefficients of \(-F\) are nonnegative.  This closes
the corner.  The complement of the local box is partitioned at
\[
 u=\frac13,\quad R=\frac12,\quad a=\frac16,\quad
 q=\frac13,\quad h=1.
\]
Twenty-nine of its rational boxes close by exact conditional subdivision,
at depth at most six.  One exceptional box has \(-F\geq0\), with its sole
zero on the weighted fan boundary.  In the other, feasibility can be an
equality only at
\(R=0,a=1/6,q=1/3,h=5\); there
\[
 F=-\frac54(3u-2),\qquad K=\frac{694-387u}{36},
\]
so $F\geq0$ again gives $K>0$.  Consequently the full infinity-face
implication $S(1)\geq0\Rightarrow S_1,J\geq0$ is exact.  The transverse
problem can now be sharpened considerably.  Write $F=S(1)$ and $H=S_1$.
In the $U^2\leq V$ chart, $H=h_0+h_1A_0$ is affine in $A_0$, with
\[
 h_1=-24Tp(D-1)(D+1)(3DW+1)(5DW+1)(6DW+1)^2(7DW+1)\geq0,
\]
whereas $F=f_0+f_1A_0+f_2A_0^2$ is convex, since
\[
 \frac{\partial^2F}{\partial A_0^2}
 =96T^2p^2(D-1)^2(D+1)^2(3DW+1)^2
 (5DW+1)(6DW+1)^3(7DW+1)\geq0.
\]
At the root $A_*=-h_0/h_1$, clear the denominator and put
\[
 \Phi=h_1^2F(A_*)=f_2h_0^2-f_1h_0h_1+f_0h_1^2.
\]
The two ordinary blow-up charts
\[
 D=R,\quad W=Rx,\qquad\hbox{and}\qquad D=Rx,\quad W=R
\]
have only nonnegative Bernstein coefficients in $-\Phi$. The arrays have
shapes $43\mathbin\times18\mathbin\times6\mathbin\times4\mathbin\times3$
and $43\mathbin\times26\mathbin\times6\mathbin\times4\mathbin\times3$.
An exact zero-face audit shows $\Phi<0$ whenever
$0<D<1$, $W>0$, $T,p>0$, including $q=0$. In the first blow-up the
maximal zero faces are
\[
R=0,T=p=1,q=0;\qquad R=1;\qquad T=0;\qquad p=0.
\]
In the second, replace $R=1$ by $R=x=1$. Every one is excluded by
these finite positive parameters. Direct Bernstein checking also gives
$F(0)<0$ on the same domain: its maximal zero faces are
\[
D=0,T=1;\quad D=q=0;\quad W=T=0;\quad T=q=0.
\]
Since $h_1>0$, convexity between $A_0=0$ and the root $A_*=-h_0/h_1$
therefore proves the strict implication
\[
                   F\geq0\quad\Longrightarrow\quad H>0.
\]
This statement includes $q=0$ and $q=1$; it is asserted on finite positive
parameters rather than every degenerate compactification face. The full
Jacobian certificate for this chart is given below.

For completeness, the Newton analysis leading to the remaining $J$ charts
is as follows.  The first singular stratum is
\[
 W=T=0,\qquad p=q=1,\qquad A_0=\frac13,
\]
with $D\in[0,1]$ free.  For $P=1-p$, $Q=1-q$, and
$C=A_0-1/3$, its first normal terms are
\[
\begin{split}
 F&=-5D^2(7D^2+8)T+O_2,\\
 H&=-(17D^2-2)T+O_2,\\
 J&=\frac13T(D^2-1)(3D^2+2)(258D^4-D^2-32)+O_2.
\end{split}
\]
Use the six weighted variables
\[
 D^2,\qquad W^2,\qquad T,\qquad 1-p,\qquad
 1-q,\qquad |A_0-1/3|.
\]
All six lower-$A_0$ charts and the upper $D^2$-dominant chart have
$-F\geq0$ coefficientwise.  In each of the three singular upper charts
where $P$, $Q$, or $C$ is dominant, put $W=R^2v$.  At the repeated
$t=v=0$ face the tangents have the form
\[
 F=a t-bv^2,\qquad H=h t-ev^2,\qquad J=j t+gv^2,
 \qquad b,e,g\geq0.
\]
The slope determinant $hb-ea$ factors with a nonnegative sign on the unit
cube.  For the Jacobian, $a>0$ implies $j\geq0$ by exact conditional
dyadic certificates: four cuts in the $P$ chart, four in the $Q$ chart,
and a depth-twelve tree in the $C$ chart with at most nine unresolved
boxes at any level.  Hence all three repeated equality faces themselves
are exact.

The full $P$-dominant radial extension is now exact as well.  In its
secondary $t$-dominant chart use
\[
 D=rd,\quad W=r^2\sigma u,\quad T=r^2\sigma^2,\quad
 p=1-r^2,\quad q=1-r^2\bar q,\quad
 A_0=\frac13+\frac23r^2c,
\]
and resolve the remaining common corner by
$r=R^2\rho$, $\sigma=R$.  After removing positive monomial factors,
the condition polynomial $F$ has multidegree
$(75,11,2,2,33,7)$ and $799$ terms, while $J$ has multidegree
$(106,16,1,3,48,8)$ and $4199$ terms.  A guided midpoint tree closes at
depth eleven with fourteen leaves: six have $J\geq0$ coefficientwise and
eight have $-F\geq0$ coefficientwise.  This is an exact integral
Bernstein computation, not a floating diagnostic.  To keep memory bounded,
the last variable is streamed in eight $F$ slices and nine $J$ slices; the
remaining slice shapes are respectively
$76\mathbin\times12\mathbin\times3\mathbin\times3\mathbin\times34$
and
$107\mathbin\times17\mathbin\times2\mathbin\times4\mathbin\times49$.
Dyadic subdivision is replayed by integral de Casteljau additions, so no
dense six-dimensional rational array is formed.

The analogous $C$-dominant radial chart is now exact.  Keep the same
$D,W,T$ substitutions, but put
\[
 p=1-r^2\bar p,\qquad q=1-r^2\bar q,\qquad
 A_0=\frac13+\frac23r^2,
 \qquad r=R^2\rho,\quad\sigma=R.
\]
In the order $(R,d,\bar p,\bar q,\rho,u)$, the condition again has
multidegree $(75,11,2,2,33,7)$ and $799$ terms.  The Jacobian has
multidegree $(106,16,4,1,48,8)$ and $4199$ terms.  Its exact midpoint
certificate has $75$ leaves at depth at most $19$: $44$ certify
$J\geq0$ and $31$ certify $-F\geq0$.  The terminal ratio $u$ needs only
the split at $1/2$.  Leaves with no $u$ cut are checked on the original
interval; the others are checked after the appropriate half-interval
substitution.  Each restricted polynomial is streamed in its terminal
Bernstein slices, and only five-dimensional integer arrays are subdivided.

For the corresponding $Q$ chart put
\[
 p=1-r^2\bar p,\qquad q=1-r^2,\qquad
 A_0=\frac13+\frac23r^2c,
 \qquad r=R^2\rho,\quad\sigma=R.
\]
The condition has the same multidegree and term count, while $J$ has
multidegree $(106,16,4,3,48,8)$ and $4199$ terms in
$(R,d,\bar p,c,\rho,u)$.  Direct subdivision reveals an additional
equality corner
\[
 R=\rho=1,\qquad d=u=\bar p=0,\qquad c=\frac14.
\]
It is resolved using the exact identity
\[
 K=J-4F,\qquad
 M=-\left.\frac{\partial^2K}{\partial d\,\partial u}
       \right|_{R=\rho=1,\ d=u=0},\qquad
 J=4F+M\left(d-\frac u2\right)^2+E.
\]
Here $M$ depends only on $(\bar p,c)$ and has $20$ terms; $E$ has
$4770$ terms and the same multidegree as $J$.  The squared term removes
the adverse mixed tangent.  At the equality corner,
\[
 M=64,\qquad
 K=480(1-R)+192(1-\rho)+96\bar p
   +\frac{640}{3}\left(c-\frac14\right)^2
   +96d^2-64du+32u^2+\cdots.
\]
The exact midpoint certificate has $38$ leaves at depth at most $35$:
$16$ use $J\geq0$, $11$ use $K\geq0$, $9$ use $-F\geq0$, and $2$
use both $M\geq0$ and $E\geq0$.  All integral Bernstein slices pass.
Thus both the $Q$- and $C$-dominant full radial implications are exact.
For a condition-negative leaf, the checker also enumerates all cube faces
on which every active Bernstein coefficient vanishes.  The Jacobian
restricted to each such face is exactly constant and nonnegative (in fact
zero except at one $C$-chart point, where it is $64/3$).  Consequently
the certificates include the equality case $F=0$ as well as $F>0$.

The full upper $W$-dominant chart is exact as well.  Use
\begin{gather*}
 D=rd,\quad W=r,\quad T=r^2t,\\
 p=1-r^2\bar p,\quad q=1-r^2\bar q,\quad
 A_0=\frac13+\frac23r^2c.
\end{gather*}
After division by $r^2$, their multidegrees and support sizes are
\[
\begin{array}{c|c|r}
 &\text{multidegree}&\text{terms}\\\hline
 F&(28,11,3,2,2,2)&799\\
 J&(40,16,3,4,1,3)&4199
\end{array}
\]
Their radial faces are $-8$ and $8(8t+3)/3$.  In particular the condition
is strictly negative on $0\leq r\leq1/2$.  The full chart has an exact
midpoint certificate with $131$ leaves at depth at most $23$: $70$ use
$-F\geq0$ and $61$ use $J\geq0$.  All condition equality faces pass the
same exact Jacobian audit.

The opposite infinity chart $V\leq U^2$ now also has an exact full radial
Jacobian certificate.  In the original coordinates $(U,W,T,p,q,A_0)$,
$F$ has multidegree $(11,3,3,2,2,2)$ and $356$ terms; $J$ has multidegree
$(16,5,3,4,1,3)$ and $1094$ terms.  Its complete midpoint tree has $393$
leaves at depth at most $25$, comprising $218$ condition-negative leaves
and $175$ Jacobian-positive leaves.  All slices and condition equality
faces have passed exact verification.

The slope implication on this opposite chart is exact for admissible
finite interior parameters.  Write $H=h_0+h_1A_0$ and
$F=f_0+f_1A_0+f_2A_0^2$.  Then
\[
\begin{split}
h_1&=-24TWp(3U+1)(5U+1)(6U+1)^2(7U+1)(U^2W-1),\\
2f_2&=96T^2Wp^2(3U+1)^2(5U+1)(6U+1)^3(7U+1)(U^2W-1)^2.
\end{split}
\]
Both have the required nonnegative sign.  The endpoint $-F(0)$ has
nonnegative Bernstein coefficients.  At the $H=0$ root, the cleared
value factors as
\[
-h_1^2F(-h_0/h_1)
=-48T^2Wp^2(3U+1)^2(5U+1)(6U+1)^3(7U+1)(U^2W-1)^2G,
\]
where $G$ is affine in $p$.  Its $p=1$ endpoint satisfies $G\leq0$
coefficientwise.  At the feasible lower endpoint $p=-h_0/(h_1/p)$,
the cleared value reduces to a four-variable polynomial.  Blow up
$(T,q)=(0,0)$ by $q=Tz$ and $T=qz$.  The resulting exact root trees
have $9$ and $47$ leaves, respectively, at depths at most $8$ and $11$.
The $A_0=1$ endpoint has exact conditional trees with $991$ and $1731$
leaves, at depths at most $27$ and $23$.  All their zero faces lie on
excluded degenerate parameter faces; the root value is strictly negative there, and convexity proves
$F\geq0\Rightarrow H>0$. The $G(p=1)$ zero faces are only
$W=0,q=0$, $W=0,q=1$, and $T=q=0$. The $F(0)$ and root-tree
zero-face audits likewise exclude no face defined by $q=0$ alone.
Thus the strict implication includes $q=0$ as needed for stationary
orientation.  This
statement does not assert the stronger implication on every face of the
closed compactification.

The full upper $T$-dominant chart is now exact too.  Set
\begin{gather*}
D=rd,\quad W=rw,\quad T=r^2,\\
p=1-r^2\bar p,\quad q=1-r^2\bar q,\quad
A_0=\frac13+\frac23r^2c.
\end{gather*}
After division by $r^2$, $F$ and $J$ have multidegrees
$(28,11,7,2,2,2)$ and $(40,16,8,4,1,3)$, with $799$ and $4199$ terms.
The equality corner is
$r=1,d=w=\bar p=0,\bar q=1,c=1/4$.  Using the same $20$-term
coefficient $M(\bar p,c)$ as in the $Q$ certificate gives
\[
J=4F+M\left(d-\frac w2\right)^2+E.
\]
The exact tree has $49$ leaves at depth at most $31$: $15$ use $J$,
$17$ use $J-4F$, $16$ use $-F$, and one uses both $M$ and $E$.
All slices and equality faces pass, including one nonconstant face on
which $J=48(3\bar q+4)$ is nonnegative.

A direct certificate now covers the entire $U^2\leq V$ chart, so the
secondary $P,Q,C$ complements need no additional covering argument.
In $(D,W,T,p,q,A_0)$, $F$ has multidegree $(11,7,3,2,2,2)$ and $356$
terms, while $J$ has multidegree $(16,8,3,4,1,3)$ and $1094$ terms.
Put
\[
\begin{split}
K&=J-4F,\qquad
M=-K_{DW}\big|_{T=1,q=0,D=W=0},\\
J&=4F+M(D-W/2)^2+E,\\
M&=-16A_0p(2220A_0^2p^3-1190A_0p^3-1170A_0p+161p+456).
\end{split}
\]
The complete tree has $95$ leaves at depth at most $29$: $49$ use
$-F\geq0$, $35$ use $K\geq0$, $9$ use $J\geq0$, and $2$ use
$M,E\geq0$. All integer slices and condition-zero-face Jacobian audits
pass. Together with the opposite chart this proves $F\geq0\Rightarrow
J\geq0$ throughout both compact cubes.

It remains to orient the stationary root; $S(1)\geq0$ must not be
assumed without proof. The exact cleared polynomials satisfy
\[
                     F(q)=F(0)+q\,C_qH(q),
\]
where $H$ is affine in $q$, $C_q$ is independent of $q$, and
\[
C_q=\begin{cases}
(1-T)(3U+1)(13U^3W+3U^2W+12U+2),&V\leq U^2,\\
(1-T)(3DW+1)(13D^3W+3D^2+12DW+2),&U^2\leq V.
\end{cases}
\]
Both factors are strictly positive for $0<T<1$. Moreover $F(0)$ is a
positive multiple of $S_0$. Suppose $S=0$ at $0<B_0\leq1$ and
$S_1\leq0$. Then $S_0=-B_0S_1\geq0$, so the strict slope implication
at $q=0$ gives $H(0)>0$. If $S_1=0$, the displayed identity at the
actual $q$ gives $F(q)=F(0)\geq0$ with $H(q)=0$, a contradiction.
If $S_1<0$, affine dependence gives a root of $H$ between $0$ and the
actual $q$; there again $F=F(0)\geq0$, contradicting strictness.
Consequently $S_1>0$, and
\[
 S(1)=(1-B_0)S_1\geq0,\qquad
 R=J/S_1\geq0,\qquad \mathcal G_4=(n+3)PR\geq0.
\]

\begin{proposition}[Feasible positive interior at level four]
For every finite $n\geq4$, every feasible positive stationary point of
\eqref{eq:k4-top-quadratic} satisfying the reduced third-moment child
bounds has nonnegative parent gap. The proof is an exact algebraic
certificate; it has not yet been formalized in Lean.
\end{proposition}

This closes the interior polynomial obligation. A claim for the entire
uniform renewal cone still requires auditing the reduction of its general
constrained boundary minima to the special slack faces discussed below.
The identities and certificate generator are in
\path{python/check_k4_uniform_interior.py}.  The $P$ checker is
\path{python/check_k4_p_radial_bernstein_tree.cpp}; the shared streamed
checker is \path{python/check_k4_radial_bernstein_tree.cpp}, with the
complete midpoint trees stored in \path{python/certificates/}.


There is now an exact result at the fully sharp corner of the reduced cone.
Put \(a=rA\), \(b=B\), \(c=rC\), \(d=D\), and write \(e,f,g,j\) for the
aggregated return and terminal moments of orders four and five.  The compact
recurrence continues as
\begin{align*}
 M_4&=4TM_3+6aM_2+4cM_1+e+f,\\
 M_5&=5TM_4+10aM_3+10cM_2+5eM_1+g+j.
\end{align*}
The lower hierarchy and \(T\geq ru\) imply the seven reduced inequalities
\begin{align*}
 nTc&\geq(n+1)a^2,&
 (n+1)ae&\geq(n+2)c^2,&
 (n+2)cg&\geq(n+3)e^2,\\
 (n-1)b&\geq nv^2,&
 nvd&\geq(n+1)b^2,&
 (n+1)bf&\geq(n+2)d^2,&
 (n+2)dj&\geq(n+3)f^2.
\end{align*}

\begin{proposition}[Fully sharp \(k=4\) renewal corner]
\label{prop:k4-fully-sharp-corner}
For every \(n\geq4\), if equality holds in all seven displayed reduced
inequalities, then
\[
 (n+3)M_3M_5-(n+4)M_4^2\geq0.
\]
\end{proposition}

The exact certificate splits at the Erlang scale \(v=(n-1)T\).  On the chart
\(T\geq v/(n-1)\), writing \(T=(v+h)/(n-1)\) produces 608 nonnegative power
coefficients.  On the other chart, write \(v=(n-1)T+h\).  For \(n\geq8\),
every coefficient is nonnegative except the middle term in
\[
 A a^6-BTha^5+CT^2h^2a^4.
\]
This quadratic is nonnegative because
\[
 4AC-B^2=
 4n^2(n-1)^{11}(n+1)^2(n+2)
 \bigl(3n^5+29n^4+111n^3+207n^2+182n+132\bigr)>0.
\]
The orders \(n=5,6,7\) require respectively four, four, and two elementary
AM--GM pairings.  For \(n=4\), the regions \(h\leq T\), \(h\geq T\) with
\(a\leq Th\), and \(h\geq T\) with \(a\geq Th\) have nonnegative
two-variable Bernstein coefficients.  These checks are exact over
\(\mathbb Q\) in \path{python/check_k4_uniform_sharp_corner.py}.

The first full one-slack extension of this proposition is also available.
Write
\[
 c=\frac{n+1}{n}\frac{a^2}{T}+s,\qquad s\geq0,
\]
and continue to impose equality in the other six reduced inequalities.

\begin{proposition}[Return-third-moment slack face]
\label{prop:k4-return-third-slack}
For every \(n\geq4\), the parent gap remains nonnegative for every \(s\geq0\)
under these assumptions.
\end{proposition}

Here is a compact description of the exact certificate.  After clearing a
positive denominator, the gap is a quartic
\[
 p_0+p_1s+p_2s^2+p_3s^3+p_4s^4.
\]
The constant term is the fully sharp gap above, and
\[
 p_4=T^4n^4(n-1)^6(n+2)>0.
\]
After positive factors are removed, \(p_3\) is a quadratic in \(a\) with
positive leading coefficient and nonpositive discriminant for \(n\geq4\).
The coefficient \(p_2\) has one AM--GM block on each of the two scale charts.
On the upper chart \(T=(v+h)/(n-1)\), every coefficient of \(p_1\) is
nonnegative.  On the lower chart \(v=(n-1)T+h\), \(p_1\) is coefficientwise
nonnegative for \(n=5\); for every \(n\geq6\), its sole negative monomial is a
multiple of
\[
 T^2ha^5s.
\]
It is absorbed by one AM--GM pairing using only one eighth of the sharp
\(a^8\) coefficient and one half of a positive \(T^4h^2a^2s^2\) coefficient.
The case \(n=4\) has two negative linear monomials; two pairings using one
sixteenth of their neighboring coefficients leave the three earlier
Bernstein charts nonnegative.  The complete rational-arithmetic check is in
\path{python/check_k4_uniform_c_slack.py}.

These two propositions still do not prove the uniform renewal lemma.  Allowing
slack in the terminal second- and third-moment inequalities (the inequalities
defining \(b\) and \(d\)) produces two boundary faces and a possible feasible
interior minimum of the top \((e,f)\)-quadratic.  The common equality boundary,
the entire return-\(c\)-slack face, and the entire terminal-\(d\)-slack face are
now settled in the stated sharp-face settings. The feasible positive
interior minimum is also settled above; the reduction of every general
constrained boundary minimum to these special faces still requires an audit.

The terminal-\(d\)-slack face reduces to a small four-variable positivity
problem whose exact certificate is now complete.  Keeping the other six
bounds sharp, write
\[
 a=zT^2,\qquad
 d=\frac{n+1}{n}\frac{b^2}{v}(1+r),\qquad z,r\geq0.
\]
After homogeneity removes \(v\), the variables are \(n-4,T/v,z,r\).
The only difficult locus found by adaptive Bernstein subdivision is
\[
 (n-1)T=v,\qquad z=r=0.
\]
If \(h\) denotes the signed deviation from this Erlang scale, the gap has
positive first-order term
\[
 \frac{n^2(n+1)^2(n+2)(n+3)}{(n-1)^7}(z+6r),
\]
while the scale deviation begins at order \(h^2\).  This leads to the weighted
blow-up \(z=h^2Z,\ r=h^2R\).  Three projective charts, according as \(h^2\),
\(z\), or \(r\) is largest, remove the local singularity.  Numerical
Bernstein subdivision then certifies every compact chart, but unstructured
exact subdivision accumulates at equality boundaries.  Weighted blow-ups,
rather than deeper rectangular subdivision, resolve these accumulations.
The reproducible exploration is in
\path{python/probe_k4_uniform_d_slack.py}; the completed exact charts are
checked by \path{python/check_k4_uniform_d_slack.py}.

The two tail normalizations have been made explicit.  In the \(T\)-dominant
chart set \(T=1\), \(v=(n-1)q\), \(a=z\), and use the actual slack
\(D=s/T^3\).  The entire region
\[
        0\leq q\leq1,\qquad 0\leq D\leq1
\]
is now certified exactly for every \(n\geq4\).  Near the return-only boundary
the correct ratio is \(R=D/q^3\).  The chart \(D=q^3R\) is coefficientwise
nonnegative in both the bounded-\(z\) and reciprocal-\(z\) charts.  In the
complementary chart put \(D=h^3\), \(q=hQ\).  Its corner has positive boundary
orders \(h^5\) and \(Q^4\), so the two charts
\[
 (h,Q)=(t^4H,t^5),\qquad (h,Q)=(t^4,t^5R)
\]
are polynomial.  The first is again coefficientwise nonnegative.  In the
second, the degree drop at \(n=\infty,R=0\) is resolved by
\(x=1/(n-3)\) and the two ordinary charts \(R=xW\), \(x=RW\).  The only
remaining projective corner there has coordinates \(U=1/z,W\), whose leading
orders are \(U^4\) and \(W^4\); its two ordinary blow-up charts close after
one midpoint split.  All these assertions are exact integer
power-to-Bernstein computations in
\path{python/check_k4_uniform_d_slack.py}.

It remains on this \(T\)-normalized half to certify the large-slack tail
\(D\geq1\).  Its natural weighted normalization is
\(H=D^{-1/3}\), \(Z=zH^2\), not the ordinary ratio \(z/D\).  Numerical
Bernstein subdivision finds no negative value in either the bounded-\(Z\) or
reciprocal-\(Z\) chart.  The singular faces \(H=0\), \(q=0\), and
\(U=1/Z=0\) are coefficientwise nonnegative.  At the remaining corner
\(H=Z=0\), the boundary orders are \(H\) and \(Z^2\).  Both weighted charts
\[
 (H,Z)=(t^2A,t),\qquad (H,Z)=(t^2,tA)
\]
have coefficientwise nonnegative faces at \(t=0\) and \(q=0\).  The formerly
missing \(A=0\) face of the second chart is the original face \(z=0\), and it
can also be closed exactly.  After removing the positive factor
\((N+3)^8(N+4)^2\), where \(N=n-4\), its gap is a quartic in \(D\).  Expanding
instead in \(E=D-1\), all five \(E\)-coefficients are nonnegative.  The only
corner not visible in their initial Bernstein arrays is
\(N=\infty,q=0\); with \(x=1/(N+1)\), the two ordinary charts \(x=qR\) and
\(q=xR\) certify every coefficient (only the linear \(E\)-coefficient needs
a dyadic split, into three boxes).  This gives a self-contained exact
certificate for \(z=0,D\geq1\) in
\path{python/check_k4_uniform_d_slack_z0.py}.

Their common residual degree staircase at
\(n=\infty,t=0\) is governed by \(1/n\sim t^3\).  Put \(x=1/(n-3)\) and use
the two charts \(x=r^3,t=rB\) and \(t=r,x=r^3B\).  All resulting powers of
\(r\) are multiples of three.  With \(y=r^3\), the lower Newton monomials at
the next corner are
\[
 q^7,\qquad yq^4,\qquad y^2q^2,\qquad y^3.
\]
The two Newton slopes give three toric regions \(y\leq q^3\),
\(q^3\leq y\leq q^2\), and \(y\geq q^2\).  The corresponding monomial charts
are now explicit.  All of their auxiliary \(B=0\) faces are coefficientwise
nonnegative; so are all \(A=0\) faces except the one already supplied by the
\(z=0\) certificate above.  These exact face checks are collected in
\path{python/check_k4_uniform_d_slack_large_faces.py}.  Thus the remaining
work initially consisted of the interiors of twelve toric charts, with all accumulating
boundary faces already proved.  Their unsplit Bernstein arrays still contain
negative coefficients, so a face-aware decomposition or controlled exact
subdivision is required; direct rectangular subdivision produces an
impractically large certificate.

One of those twelve charts is now completely closed.  In the
\(H\)-dominant, \(t\)-dominant chart with \(y\leq q^3\), all
negative Bernstein coefficients accumulate at \(q=A=0,\ R=B=1\).  Its next
Newton terms are
\[
 A^4,\qquad A^3q,\qquad A^2q^2,\qquad Aq^3,\qquad q^6.
\]
The polynomial has the manifestly positive factor
\[
              (1+2BR^3q^3)^8.
\]
After removing it, the three regions \(A\leq q^3\),
\(q^3\leq A\leq q\), and \(q\leq A\) have exact nonnegative Bernstein
coefficients.  The latter two would create large dense four-dimensional
tensors, so the conversion is performed one \(B\)-slice at a time.  The
complete low-memory certificate is
\path{python/check_k4_uniform_d_slack_large_ht_first.py}.

Two further toric interiors close after the common positive factor
\((1+2w^6)^8\) is removed.  The \(H\)-dominant, \(x\)-dominant chart with
\(y\geq q^2\) requires five terminal dyadic boxes and depth at most six; the
corresponding \(Z\)-dominant, \(x\)-dominant chart requires nine boxes and
depth at most ten.  Their exact checkers are
\path{python/check_k4_uniform_d_slack_large_hx_last.py} and
\path{python/check_k4_uniform_d_slack_large_zx_last.py}.  Thus three of the
original twelve interiors are complete at this stage.

The \(Z\)-dominant, \(t\)-dominant chart with \(y\leq q^3\) is also
complete.  After removing
\((1+2BR^3q^3)^8\), two dyadic splits in \(A\) certify the whole region
\(1/4\leq A\leq1\); see
\path{python/check_k4_uniform_d_slack_large_zt_first_large_a.py}.  At \(A=0\)
the reduced polynomial factors as
\[
 q^2R^9B^{10}(1+3BR^3q^3)(1+4BR^3q^3)
 (1+5BR^3q^3)(1+6BR^3q^3)^2.
\]
The nondominated exponent vectors in the order \((A,q,R,B)\) are
\[
\begin{split}
 &(0,2,9,10),(1,1,6,8),(2,0,3,6),(4,1,3,5),\\
 &(5,0,0,3),(7,1,0,2),(9,2,0,1),(11,3,0,0).
\end{split}
\]
The small-\(A\) tail is a genuine multiscale corner.  Its first comparison
is \(A\sim qR^3B^2\), and the reason for this scale is especially
transparent: the first three Newton monomials, including their signs, are
exactly
\[
 q^2R^9B^{10}-2AqR^6B^8+A^2R^3B^6
   =R^3B^6(A-qR^3B^2)^2.
\]
The face factorization, Newton exponents, and leading square are checked
exactly in
\path{python/check_k4_uniform_d_slack_large_zt_first_newton.py}.

These scales lead to a short complete cover.  Below the first surface, put
\(A=qR^3B^2C\) and \(U=1-C\).  The only residual equality corner has Newton
orders
\[
                 U^2,\qquad q^3U,\qquad q^4.
\]
The three associated charts \(U\leq q^3\),
\(q^3\leq U\leq q\), and \(q\leq U\) have exact nonnegative Bernstein
certificates; only a depth-twelve dyadic refinement of the first chart is
needed.  This is checked in
\path{python/check_k4_uniform_d_slack_large_zt_first_lower.py}.

For \(qR^3B^2\leq A\leq RB\), write \(A=RBX\).  The condition becomes
\(qR^2B\leq X\), and four ordinary projective charts give direct streamed
Bernstein certificates; see
\path{python/check_k4_uniform_d_slack_large_zt_first_middle.py}.  Finally,
\(RB\leq A\) is covered by the two charts \(R\leq A\) and
\(A\leq R, RB\leq A\), checked in
\path{python/check_k4_uniform_d_slack_large_zt_first_upper.py}.  Hence four
of the original twelve toric interiors are now complete and eight remain.

The adjacent \(H\)-dominant, \(t\)-dominant middle chart is likewise
complete.  Removing \((1+2BQ^6C^3)^8\) leaves the boundary factor
\[
                   A^2B^6(A-B^2C)^2\qquad (Q=0).
\]
On the side \(A\leq B^2C\), the substitution \(A=B^2CX\), \(U=1-X\)
has only the transverse Newton orders \(U^2\) and \(Q^6\); the two charts
at \(U=Q^3\) have nonnegative exact Bernstein coefficients.  Three ordinary
projective charts cover \(B^2C\leq A\).  The streamed certificate is
\path{python/check_k4_uniform_d_slack_large_ht_middle.py}.  Thus five of the
twelve toric interiors are complete and seven remain.

The final \(H\)-dominant, \(t\)-dominant chart is complete as well.
After removing \((1+2Bw^6)^8\), its boundary factor is
\[
                   A^2B^6(AC^2-B^2)^2\qquad (w=0).
\]
Four monomial projective charts cover the two sides of \(AC^2=B^2\).
One chart requires a single dyadic split of its ratio variable; the five
resulting streamed Bernstein checks are all nonnegative.  See
\path{python/check_k4_uniform_d_slack_large_ht_last.py}.  Thus six of the
twelve toric interiors are complete and six remain.  The exploratory transformations are in
\path{python/probe_k4_uniform_d_slack_large_toric.py}.  In the
\(Z\)-dominant, \(x\)-dominant chart with \(y\leq q^3\), the Newton boundary
starts with
\[
             R^3(AB^3-qR^3)^2+A^5B^{12}.
\]
Five monomial charts give a complete exact certificate on the side
\(qR^3\leq AB^3\); see
\path{python/check_k4_uniform_d_slack_large_zx_first_above.py}.  On the
opposite side put
\[
             B=R(1-U),\qquad A=q(1-V).
\]
After the common nonnegative factor \(q^2R^9\) is removed, its leading core
is
\[
 (3U+V)^2-12(3U+V)q^3R^3+8q^4R^6.
\]
Thus the transverse scale is
\(\max\{U,V\}\asymp q^2R^3\).  Split into \(V\leq U\) and \(U\leq V\),
and write \(D=\max\{U,V\}\), with the smaller deficit equal to \(DZ\).  For
\(D\leq q^2R^3\), put \(D=q^2R^3W\).  The resulting compact chart has a
direct streamed Bernstein certificate on the three \(R\)-intervals
\[
 [0,1/2],\qquad[1/2,3/4],\qquad[3/4,1];
\]
only the last needs one midpoint split in \(W\).

For \(q^2R^3\leq D\), five regular monomial charts give a complete cover:
\[
\begin{array}{c|c}
 R\leq D & R=DS,\\
 R^2\leq D\leq R & R=ST,\quad D=S^2T,\\
 R^3\leq D\leq R^2 & R=TL,\quad D=T^3L^2,\\
 qR^3\leq D\leq R^3 & D=R^3L,\quad q=LM,\\
 q^2R^3\leq D\leq qR^3 & D=R^3M^2N,\quad q=MN.
\end{array}
\]
After division by the corresponding leading monomial, all ten charts (five
for each deficit ordering) have nonnegative exact Bernstein coefficients.
The four-dimensional conversions are streamed one final-axis slice at a
time, and the five large charts are constructed and discarded sequentially.
The certificate for this side is
\path{python/check_k4_uniform_d_slack_large_zx_first_below.py}; the wrapper
\path{python/check_k4_uniform_d_slack_large_zx_first.py} checks both sides.
Hence the first \(Z\)-dominant, \(x\)-dominant interior is complete: seven of
the original twelve toric interiors are now closed and five remain.

The corresponding first \(H\)-dominant, \(x\)-dominant interior is now
complete as well.  The same two-sided projective cover applies after
canceling the smaller common monomials appropriate to the \(H\)-chart.  On
the side \(qR^3\leq AB^3\), the five streamed polynomials have only \(751\)
terms each.  On the opposite side, the two compact charts and ten
large-deficit charts have \(15948\) terms each; their maximal degrees are
strictly smaller than those of the \(Z\)-dominant certificate.  Every
streamed Bernstein coefficient is nonnegative, with the same single
midpoint split of \(W\) on the top compact \(R\)-interval.  The complete
low-memory checker is
\path{python/check_k4_uniform_d_slack_large_hx_first.py}.  Thus eight of the
twelve toric interiors are complete and four remain.

The middle \(H\)-dominant, \(x\)-dominant interior is now complete too.
After removing \((1+2Q^6C^3)^8\), every exponent of \(Q\) and \(B\) is a
multiple of three.  Thus the substitutions
\[
                         P=Q^3,\qquad K=B^3
\]
reduce the degrees \(267\) and \(183\) in the first attempted certificate to
\(89\) and \(61\).  At \(P=0\) the compressed polynomial has the exact face
\[
                         A^2(AK-C)^2.
\]
On \(C\leq AK\), put \(C=AKX\), \(U=1-X\); the two ordinary charts
\(U=PS\) and \(P=US\) have nonnegative streamed Bernstein coefficients.
The opposite side is covered without denominators by
\[
 A\leq C:\quad A=CX,\qquad
 C\leq A:\quad C=AY,\quad K=YX.
\]
In the first chart one orders the deficits \(1-X\) and \(1-K\), then splits
their maximum against \(P\).  In the second chart the first residual scale
is again \(U\sim P\).  Its only further boundary accumulation has face
\(S^2\) at \(Y=0\), hence the weighted split \(Y\sim S^2\).  The remaining
Newton boundary is
\[
                         R^2,\qquad T^5R,\qquad T^6,
\]
so a final split at \(R=T^3\), followed by a square compression of the radial
variable, regularizes the corner.  Three rational intervals and one shallow
terminal subdivision close the two far boxes.  The complete low-memory
certificate is
\path{python/check_k4_uniform_d_slack_large_hx_middle.py}.
Consequently nine of the twelve toric interiors are complete.  The three
remaining ones are the middle \(Z\)-\(x\) chart and the middle and final
\(Z\)-\(t\) charts.

The middle \(Z\)-dominant, \(x\)-dominant interior is now complete as well.
The same cube compression gives a (751)-term polynomial of degree
\((38,53,13,8)\) in \((P,C,A,K)\), but its face is now
\[
                              (AK-C)^2.
\]
The side \(C\leq AK\) is again covered by \(C=AKX\), \(U=1-X\), and the
two charts \(U=PS\), \(P=US\).  On the opposite side, the charts \(A=CX\)
and \(C=AY, K=YX\) give a polynomial cover of \(A\leq C\) and \(C\leq A\),
respectively.  In the first chart, order (1-X) and (1-K) and compare
their maximum with \(P\).  Two of the four resulting charts are direct.
The other two have the Newton boundary
\[
                         C^4,\qquad C^3W,\qquad W^2,
\]
which is resolved by
\[
 W=C^3T,\qquad (C,W)=(LT,LT^3),\qquad C=WT.
\]
The low chart uses four rational \(T\)-intervals, the middle chart two
dyadic intervals, and the high chart is direct.

For \(C\leq A\), the successive scales are \(U\sim P\), \(Y\sim S^2\),
and \(R\sim T^3\), exactly as in the \(H\)-dominant calculation.  On the
low \(R\)-chart, \(W\geq1/2\) is direct after three rational intervals.
After rescaling the remaining half to a unit box, its only unresolved face
has Newton monomials
\[
                         A^4,\qquad A^3W,\qquad W^2,
\]
and the three charts
\[
 W=A^3T,\qquad (A,W)=(LT,LT^3),\qquad A=WT
\]
all have nonnegative streamed Bernstein coefficients.  The opposite
\(T^3\leq R\) chart closes on two dyadic radial intervals.  The complete
low-memory checker, with independently selectable sections, is
\path{python/check_k4_uniform_d_slack_large_zx_middle.py}.  Consequently ten
of the twelve toric interiors are complete; only the middle and final
\(Z\)-\(t\) charts remain.

The final \(Z\)-dominant, \(t\)-dominant interior is now complete.  After
removing \((1+2Bw^6)^8\), its face at \(w=0\) is
\[
                         B^6(AC^2-B^2)^2.
\]
Four ordinary projective regions cover the two sides of this equality.  The
only additional accumulation occurs in the \(B\leq C\) region.  After the
common factor \(X^{10}\) is removed, its equality deficit has the weighted
scale \(1-Y\asymp w^3\).  The high chart \(w^3\leq1-Y\) is direct.  In the
low chart the Newton boundary is
\[
 S^2,\qquad 8X^2,\qquad -12Cw^3XS,\qquad C^2w^6.
\]
Comparing \(S\), \(X\), and \(w^3\) gives a three-way weighted fan.  The
\(S\)-dominant chart is direct, the \(X\)-dominant chart closes by a
depth-twelve streamed subdivision, and the \(w^3\)-dominant chart reduces
to the positive face \(V^2+8U^2+C^2\).  An ordinary three-way fan comparing
\(U,V,C\) resolves that face; only the \(U\)- and \(C\)-dominant charts need
depth-twelve subdivision.  The ten independently selectable checks are in
\path{python/check_k4_uniform_d_slack_large_zt_last.py}.  Consequently eleven
of the twelve toric interiors are complete; only the middle \(Z\)-\(t\)
interior remains.

The middle \(Z\)-dominant, \(t\)-dominant interior is complete as well.
After removing \((1+2BQ^6C^3)^8\), its face at \(Q=0\) is
\[
                         B^6(A-B^2C)^2.
\]
On \(A\leq B^2C\), put \(A=B^2CX\) and \(U=1-X\).  The far chart
\(Q^3\leq U\) is direct.  In the near chart \(U\leq Q^3\), the remaining
Newton boundary is
\[
                         C^4,\qquad C^3S,\qquad S^2.
\]
The associated low, middle, and high fan charts cover \(S\leq C^3\),
\(C^3\leq S\leq C\), and \(C\leq S\).  The middle chart closes on the two
dyadic halves of its radial variable, and the high chart is direct.  In the
low chart, after the cube compression \(P=Q^3\), the secondary Newton
support is
\[
                         B^2,\qquad CP^2,\qquad C^2.
\]
Eight ordinary projective charts compare these three scales.  Their
four-dimensional Bernstein conversions are streamed one final-axis slice
at a time; only terminal slices require depth-eighteen subdivision.  The
opposite side \(B^2C\leq A\) is covered by three ordinary projective charts,
all direct after their common positive monomial factors are removed.  The
complete low-memory certificate, with independently selectable sections
and fixed degree-and-support assertions, is
\path{python/check_k4_uniform_d_slack_large_zt_middle.py}.

Consequently all twelve toric interiors, together with their previously
certified boundary faces, are closed.  This completes the \(T\)-normalized
large-slack tail of the terminal-\(d\)-slack face.

In the
\(v\)-dominant chart set \(v=1\),
\(T=q/(n-1)\), \(a=zT^2\), retain the relative slack \(r\), and divide the
cleared numerator by its common factor \(q^4\).  A further ordinary
projective split according as \(1\), \(z\), or the slack is largest removes
the artificial degree drop at infinity.  The central and slack-dominant
charts are direct.  The return-dominant chart has only the corner
\(N=q=U=V=0\).  Its Newton facets in \((q,U,V)\) have primitive normals
\[
                  (1,2,1),\qquad(2,3,3),\qquad(4,5,5).
\]
After ordering \(U,V\), the side \(V\leq U\) is covered by the radial scales
\(R\asymp q^{3/2}\) and \(R\asymp q^{5/4}\), while the side \(U\leq V\) has a
homogeneous degree-eight three-way fan.  The infinity faces of these charts
are resolved by the successive Newton normals \((1,1),(2,5)\), then by a
homogeneous degree-three fan, the \((2,3)\) face, and finally the
\((1,4,3)\) face.

The last accumulation occurs in a low-\(t\), \(A\)-dominant ratio chart.  Its
\((A,R)\)-Newton boundary is
\[
                    (0,24),\qquad(1,18),\qquad(2,0),
\]
with respective leading coefficients \(1,-2Xt^4,1+X\).  The middle point
lies strictly above the facet of normal \((12,1)\).  Hence the two weighted
charts \(A=R^{12}B\) and \(A=D^{12},R=DB\) have positive leading faces
\(1+B^2(1+X)\) and \(B^{24}+1+X\), respectively, and their exact Bernstein
certificates close after four and two subdivisions.  The analogous residual
\(\rho\)-chart has boundary \((0,4),(1,3),(2,0)\) and is closed by the two
charts associated with the normal \((2,1)\).  Replacing the recurring powers
\(R^6,\rho^6,t^2\) by new unit-interval variables keeps the largest dense
certificate small.  Fixed degree-and-support assertions and the complete
exact checks are in
\path{python/check_k4_uniform_d_slack_v_tail.py}.

Together with the \(T\)-normalized certificates above, this completes the
entire terminal-\(d\)-slack boundary face under its six sharp bounds.
The feasible positive interior minimum is now certified above. Before
claiming the uniform \(k=4\) renewal lemma, the general boundary reduction
must still be justified.

The exact constrained-boundary audit keeps the lower moments free.
Normalize $M_1=1$, put $\alpha=(n+3)/(n+2)$,
$\beta=(n+2)/(n+1)$, and write
\[
 y_0=\beta c/a,\qquad f_0=\beta d/b,\qquad
 \delta_0=5/(2\alpha).
\]
For $Q=\mathcal G_4$ in the ratio coordinates $(y,f)$,
\[
 \frac{Q_y}{c}-\frac{Q_f}{d}
     =2(n+3)P\alpha(y-f+\delta_0).
\]
Thus a minimum on $y=y_0$ with $f>f_0$ must have $Q_f=0$ and
$f-y_0\leq\delta_0$. Its free-coordinate root is
\[
 f_e=\frac{2(n+4)(L+cy_0)-5(n+3)PT}
              {2\{(n+3)P\alpha-(n+4)d\}}.
\]
Likewise a minimum on $f=f_0$ with $y>y_0$ has $Q_y=0$ and
$f_0-y\geq\delta_0$, with root
\[
 y_e=\frac{2(n+4)(L+df_0)-5(n+3)P(T+1)}
              {2\{(n+3)P\alpha-(n+4)c\}}.
\]
Both denominators are strictly positive. Completing the two squares is
exact; the cleared edge-root values have respectively $967$ and $1078$
terms in $(n,T,a,b,c,d)$. Their conditional nonnegativity under root
feasibility and the displayed gradient signs remains open, as does the
general corner $y=y_0,f=f_0$ with simultaneous lower-moment slacks.
These identities and support assertions are reproducible in
\path{python/check_k4_uniform_top_boundaries.py}; they are reductions,
not positivity certificates for the missing faces.

An additive-cone shortcut does not close this gap.  If the moment triple is
split into its return-only part \(X\) and terminal-response part \(Y\), the
return-only part satisfies the desired log-convexity inequality, but \(Y\)
need not.  Already at \(n=4,T=1,v=2,a=s=0\), its corresponding gap is
\(-819136/243\).  Thus closure of the log-convexity cone under addition cannot
replace the remaining cross-term analysis.

Monotonicity or convexity in the terminal slack also fails.  For the cleared
\(T\)-normalized numerator, at \(n=4,q=1,z=3,D=1\) the first \(D\)-derivative
is \(-1804714429536\).  At \(n=6,q=1,z=11,D=122\), the second derivative is
\(-490528950000000\).  Consequently the certified face \(D=1\) cannot simply
be propagated through the large-slack tail by either argument.

One should not try to repeat the next convexity step from the \(k=3\) proof.
After setting
\[
 E=\frac{n+2}{n+1}\frac{C^2}{A},\qquad
 F=\frac{n+2}{n+1}\frac{D^2}{B},
\]
the reduced gap is not convex in \(C,D\).  For example, at
\[
 n=4,\quad x=r=u=v=1,\quad A=\frac43,\quad B=2,\quad
 C=\frac{20}{9},\quad D=5,
\]
its Hessian is
\[
 \begin{pmatrix}
  558194/15&-1975/2\\
  -1975/2&-155882/15
 \end{pmatrix},
 \qquad
 \det=-\frac{348927229057}{900}.
\]
The gap itself is nevertheless positive there. The interior stationary
lemma above is now certified, but nested convexity cannot supply the
missing general boundary reduction.  All identities and this exact witness are checked by
\path{python/probe_k4_uniform_renewal.py}.

\section{The acyclic case}

Conjecture~\ref{conj:matrix} has a short proof when the directed graph of the
off-diagonal entries of \(A\) is acyclic.  Expanding an entry of \(e^{tA}\) by
directed paths, a path with \(\ell\) off-diagonal edges contributes
\[
  t^\ell
  \int_{\Delta_\ell}
  \exp\!\left(t\sum_{j=0}^{\ell}x_j a_{i_ji_j}\right)d\mu(x),
\]
up to a positive constant.  Since the path is simple, \(\ell\leq s-1\).
After division by \(t^{s-1}\), every integrand is of the form
\[
  t^{-(s-1-\ell)}e^{ct},
\]
which is log-convex.  Positive integrals and positive sums of log-convex
functions are log-convex (by H\"older's inequality).  This recovers the
triangular result and shows precisely where cycles enter: a cyclic walk may have
arbitrarily many off-diagonal edges, so the exponent \(s-1-\ell\) is no longer
nonnegative path by path.

\section{Evidence and completion in dimension three}

The curvature in \eqref{eq:matrix-curvature} was tested numerically on random
dense and sparse nonnegative matrices through size \(6\), including extreme rate
ratios, deterministic initial and terminal vectors, directed cycles, and direct
optimization of the curvature.  No negative value was found; serial acyclic
chains approach equality.  This is evidence only, not a proof.

In dimension three, Proposition~\ref{prop:three-reduction} makes a search over
general matrices unnecessary.  Direct optimization of the three canonical
curvatures over widely separated rates and feedback parameters likewise found
no negative value.  Taylor-coefficient experiments reveal a more precise
obstruction: after the shift leading to \eqref{eq:canonical-rate-system}, the
second and third start states have many manifestly nonnegative coefficient
groups, whereas the first start has negative monomials beginning in feedback
degree two.  The full coefficient and the pointwise curvature remain
nonnegative in all tests.  Thus a proof must group unequal-rate terms (typically
into expressions involving $(x_i-x_j)^2$); positivity monomial by monomial is
already false in this reduced chart.

Proposition~\ref{prop:complex-middle} and the real-root divided-difference
argument settle the middle start for every spectrum; convolution then settles
the third start.  Proposition~\ref{prop:complex-first} now settles the first
start throughout the complex-spectrum region as well.  The degree-$60$
certificate following \eqref{eq:real-root-compact-chart} also settles the
real-spectrum first start through Perron-normalized time $10$.
Proposition~\ref{prop:real-first-normalized-complete-logconvex}, the coefficient
comparison following \eqref{eq:real-first-baseline-coefficient}, and the exact
low-degree block \eqref{eq:real-first-variance-low-block} settle its unbounded
tail.

\begin{corollary}[Complete three-state case]
\label{cor:complete-three-state}
Conjecture~\ref{conj:matrix} holds for every $3\times3$ nonnegative matrix.
Consequently, through Proposition~\ref{prop:fisher}, the full moment hierarchy
holds for every three-state phase-type law.
\end{corollary}

\paragraph{Formalization status.}
The Lean development now proves normalized complete-symmetric log-convexity
for every index.  The exact block covers
$1\leq n\leq5$, and the transition argument above covers $n\geq6$.
For the last transition-point estimate, the formal proof uses the simpler
majorant
\[
 b_n\leq\frac{n-1}{n+1}\qquad(n\geq7),
\]
deduced from $(1+1/m)^m\leq e<3$; the case $n=6$ is an exact tenth-power
comparison.  This avoids formalizing the auxiliary real-variable monotonicity
argument while proving the same bound.  The combined Lean theorem is
\texttt{normalizedComplete\_logConvex\_all}.  The subsequent weighted
Chebyshev reduction has also been formalized: all dependence on $u,v$ is
discharged.  The parameter-free sum in
\eqref{eq:real-first-baseline-coefficient} is now evaluated in Lean by expanding
its kernel in the first five falling-factorial polynomials and using the
binomial falling-factorial moments.  Consequently the formal far-tail
coefficient theorem has no unevaluated finite-sum hypothesis.

The remaining finite work has since been shortened further.  After restoring
the constant term $h_1=1+w$, the full corrected-curvature coefficients are
nonnegative directly: degrees $N=6,7,8$ have short algebraic proofs, degrees
$N=9,\ldots,14$ use six independent two-variable Bernstein certificates, and
the general comparison proves every $N\geq15$.  The certificates are isolated
in separate modules so that they can be checked with bounded peak memory.
Their assembly is
\texttt{fullCurvatureCoefficient\_nonneg\_of\_root}.

Absolute convergence and two termwise differentiations are also formalized.
The full coefficient series is identified with
\[
 t^2\{h_1h_1''-(h_1')^2\}+2h_1^2,
\]
and summed term by term.  Thus
\texttt{firstStartCurvature\_nonneg\_of\_root\_all\_time}, and equivalently
\texttt{normalizedFirstStart\_correctedCurvature3\_nonneg}, proves the
real-spectrum first start for every $t\geq0$; the former restriction
$t\geq10$ and the degree-$60$ compact-time certificate are no longer needed.

The normalization back toward the canonical cubic is now partly formalized as
well.  The file
\texttt{K3CyclicRealCanonicalBridge.lean} proves amplitude, time-scale, and
exponential-shift transport, including a version stated with actual first and
second derivatives.  From Vieta's identities and the canonical sign
conditions it proves that ordered real roots are nonnegative and that their
ratios $u=\lambda_1/\lambda_3$, $v=\lambda_2/\lambda_3$ satisfy
$\sqrt u+\sqrt v\leq1$.  Consequently
\path{canonicalOrderedRealRoots_firstStart_nonneg} is the checked
pointwise first-start inequality for every ordered real-root canonical cubic.
The remaining matrix-exponential identification is now formalized as well.
The file \texttt{K3CyclicCanonicalMatrixBridge.lean} proves the canonical
matrix cubic identity, differentiates the actual first coordinate of the
matrix exponential three times, and invokes uniqueness for the resulting
scalar cubic ODE.  Its theorem
\path{canonicalFirst_eq_orderedRealRootSeries} identifies that coordinate
with the normalized Taylor series, while
\path{canonicalFirst_orderedRealRoots_correctedCurvature3_nonneg} states the
real-spectrum certificate directly for the actual first coordinate and its
actual derivatives.

\begin{sloppypar}
The complex first start is now complete in Lean.  The lightweight identity
module \path{K3CyclicComplexFirst.lean} contains the oscillator energy,
Wronskian factorization, and reduction to
\eqref{eq:first-oscillator-target}; the certificate and analytic-tail modules
cover the compact and large-$x$ regions.  Their assembly in
\path{K3CyclicComplexFirstFull.lean} is
\path{firstH_correctedCurvature3_nonneg}.  The ODE-uniqueness bridge
\path{K3CyclicComplexCanonicalFirstBridge.lean} now identifies this normalized
function with the actual first coordinate of the canonical matrix exponential.
The theorem
\path{canonicalFirst_complexPair_correctedCurvature3_nonneg_of_vieta} in
\path{K3CyclicComplexCanonicalFirstCurvature.lean} transports the certificate
to that actual coordinate; its compact-domain hypotheses are derived internally
from the canonical sign conditions and the complex Vieta identities.
\end{sloppypar}

The coefficient inequality behind the complex middle start is also now
formalized.  The exact conjugate pairing is proved in
\path{K3CyclicComplexMiddlePairingReal.lean}.  Separate low-memory Taylor
modules prove \eqref{eq:middle-J-remainder}, including the sharp third-
derivative estimate for $(1-it)^2(1+it)^{N-2}$.  The finite Bernstein
certificate covers $4\leq N\leq24$, while
\path{K3CyclicComplexMiddleTailCoefficient.lean} combines the analytic tail,
the numerical domination for $N\geq25$, and the degenerate $\omega=0$
boundary.  Its theorem
\path{scalarCurvatureCoefficient_nonneg_all} proves every middle curvature
coefficient is nonnegative throughout the compact chart.  The EGF and Cauchy-
product modules \path{K3CyclicComplexMiddleSeries.lean} and
\path{K3CyclicComplexMiddleCurvatureSeries.lean} now sum those coefficients
and prove \path{middleRecurrenceCurvature_nonneg}, the analytic-series form of
\eqref{eq:complex-root-middle-curvature}.  That identification is now complete:
\path{K3CyclicComplexCanonicalMiddleBridge.lean} differentiates the shifted and
scaled EGF, matches its cubic ODE and initial data to the actual middle
coordinate, and transports the stronger curvature inequality.  The separate
low-memory module \path{K3CyclicComplexCanonicalMiddleChart.lean} proves the
two square identities giving the exact interval
\eqref{eq:complex-middle-A-interval}, constructs the compact chart variables
from the complex Vieta data, and concludes with
\path{canonicalMiddle_complexPair_middleCurvature3_nonneg}.  Finally,
\path{K3CyclicCanonicalThird.lean} formalizes the actual third coordinate and
the driven identities $h_3'=x_3h_2+bh_3$ and
$h_3''=x_3h_2'+bh_3'$.  The companion module
\path{K3CyclicCanonicalThirdIntegral.lean} proves the exact variation-of-
constants formula used above.  The remaining formal work is the general
convolution log-convexity lemma carrying the stronger middle inequality to the
third coordinate, followed by the final canonical-to-general matrix reduction.

The paper \emph{Log-Convex Matrix Functions} by Aujla, Rawla, and Vasudeva uses
``matrix log-convexity'' in a different sense: it studies the functional calculus
\(A\mapsto f(A)\) for positive-definite Hermitian matrices.  Its main criteria do
not apply to scalar coefficients \(u^Te^{tA}v\) with a nonsymmetric nonnegative
matrix \(A\).  Its closure of positive sums of log-convex functions agrees with
the elementary closure used in the acyclic argument above, but does not prove
Conjecture~\ref{conj:matrix}.

The paper by Butt, Pe\v{c}ari\'c, and Rehman on exponential convexity of
Petrovi\'c-type functionals is also about a different class of convexity
preservers.  It contains no renewal, phase-type, or nonsymmetric semigroup
result that applies to \eqref{eq:canonical-laplace}.

\end{document}
